% ============================================================================
%  gsd-going-further-guide.tex
%  GSD: Going Further --- Skills, MCP, and Multi-Repo Workflows
%  Compile: xelatex (3-pass) from this directory; \input@path picks up shared/
% ============================================================================
\documentclass[11pt,letterpaper,openany]{book}

\title{GSD: Going Further}
\author{Tibsfox}
\date{}

\makeatletter
\def\input@path{{../shared/}}
\makeatother

% ============================================================================
%  gsd-guides-preamble.tex
%  Shared LaTeX preamble for the GSD User Guides Suite
%  Compile with: xelatex (3-pass), NOT pdflatex
%  License: CC BY-SA 4.0 — see \cclicense macro
%
%  Load order note: this file loads xcolor + tcolorbox BEFORE sourcing
%  gsd-color-scheme.tex, and then loads hyperref AFTER color definitions
%  are live. Documents only need to % ============================================================================
%  gsd-guides-preamble.tex
%  Shared LaTeX preamble for the GSD User Guides Suite
%  Compile with: xelatex (3-pass), NOT pdflatex
%  License: CC BY-SA 4.0 — see \cclicense macro
%
%  Load order note: this file loads xcolor + tcolorbox BEFORE sourcing
%  gsd-color-scheme.tex, and then loads hyperref AFTER color definitions
%  are live. Documents only need to % ============================================================================
%  gsd-guides-preamble.tex
%  Shared LaTeX preamble for the GSD User Guides Suite
%  Compile with: xelatex (3-pass), NOT pdflatex
%  License: CC BY-SA 4.0 — see \cclicense macro
%
%  Load order note: this file loads xcolor + tcolorbox BEFORE sourcing
%  gsd-color-scheme.tex, and then loads hyperref AFTER color definitions
%  are live. Documents only need to \input{gsd-guides-preamble} in their
%  preamble (this file pulls in the color scheme transitively).
% ============================================================================

% ---------- Fonts ----------
\usepackage{fontspec}
\setmainfont{DejaVu Serif}[
  Ligatures=TeX,
  Scale=0.95,
  BoldFont={DejaVu Serif Bold},
  ItalicFont={DejaVu Serif Italic},
  BoldItalicFont={DejaVu Serif Bold Italic}
]
\newfontfamily\headingfont{DejaVu Sans}[Scale=0.95]
\newfontfamily\monofont{DejaVu Sans Mono}[Scale=0.88]
\setmonofont{DejaVu Sans Mono}[Scale=0.88]

% ---------- Page geometry ----------
\usepackage[
  letterpaper,
  top=1in,
  bottom=1in,
  left=1.25in,
  right=1.25in,
  headheight=14pt
]{geometry}

% ---------- Color + callout engine (must load BEFORE color-scheme input) ----
\usepackage{xcolor}
\usepackage[most]{tcolorbox}

% ---------- Typesetting packages ----------
\usepackage{tabularx}
\usepackage{longtable}
\usepackage{booktabs}
\usepackage{colortbl}
\usepackage{lastpage}
\usepackage{enumitem}
\usepackage{listings}
\usepackage{tikz}
\usepackage{fancyhdr}
\usepackage{titlesec}
\usepackage{parskip}
\usepackage{microtype}

% ---------- Pull in the color scheme (defines gsdnavy, gsdaccent, callouts) -
\input{gsd-color-scheme.tex}

% ---------- Hyperref (needs colors already defined) ----------
\usepackage[
  colorlinks=true,
  linkcolor=gsdnavy,
  urlcolor=gsdaccent,
  citecolor=gsdnavy,
  bookmarksnumbered=true
]{hyperref}

% ---------- Header / footer ----------
\pagestyle{fancy}
\fancyhf{}
\fancyhead[L]{\small\headingfont\@title}
\fancyhead[R]{\small\headingfont\leftmark}
\fancyfoot[L]{\small\headingfont GSD v1.35.0 — CC BY-SA 4.0}
\fancyfoot[R]{\small\headingfont Page \thepage\ of \pageref*{LastPage}}
\renewcommand{\headrulewidth}{0.4pt}
\renewcommand{\footrulewidth}{0.4pt}

\fancypagestyle{plain}{%
  \fancyhf{}
  \fancyfoot[L]{\small\headingfont GSD v1.35.0 — CC BY-SA 4.0}
  \fancyfoot[R]{\small\headingfont Page \thepage\ of \pageref*{LastPage}}
  \renewcommand{\headrulewidth}{0pt}
  \renewcommand{\footrulewidth}{0.4pt}
}

% ---------- Chapter / section styling ----------
\titleformat{\chapter}[display]
  {\headingfont\huge\bfseries\color{gsdnavy}}
  {\chaptertitlename\ \thechapter}{20pt}{\Huge}
\titleformat{\section}
  {\headingfont\Large\bfseries\color{gsdnavy}}
  {\thesection}{1em}{}
\titleformat{\subsection}
  {\headingfont\large\bfseries\color{gsdnavy}}
  {\thesubsection}{1em}{}

% ---------- Listings (code blocks) ----------
\lstdefinestyle{gsdcode}{
  basicstyle=\ttfamily\small,
  backgroundcolor=\color{gsdcodebg},
  frame=leftline,
  framerule=2pt,
  rulecolor=\color{gsdaccent},
  xleftmargin=1em,
  xrightmargin=0.5em,
  aboveskip=0.8em,
  belowskip=0.8em,
  breaklines=true,
  breakatwhitespace=true,
  showstringspaces=false,
  columns=flexible,
  keepspaces=true
}
\lstset{style=gsdcode}

% ---------- Semantic macros ----------
\newcommand{\cclicense}{%
  \textit{Released under CC BY-SA 4.0.\\
  Authored by Tibsfox. GSD itself is MIT-licensed — see gsd-build/get-shit-done.}%
}

\newcommand{\chapterheading}[1]{{\headingfont\LARGE\bfseries\color{gsdnavy}#1\par\medskip}}
\newcommand{\sectionheading}[1]{{\headingfont\Large\bfseries\color{gsdnavy}#1\par\smallskip}}
 in their
%  preamble (this file pulls in the color scheme transitively).
% ============================================================================

% ---------- Fonts ----------
\usepackage{fontspec}
\setmainfont{DejaVu Serif}[
  Ligatures=TeX,
  Scale=0.95,
  BoldFont={DejaVu Serif Bold},
  ItalicFont={DejaVu Serif Italic},
  BoldItalicFont={DejaVu Serif Bold Italic}
]
\newfontfamily\headingfont{DejaVu Sans}[Scale=0.95]
\newfontfamily\monofont{DejaVu Sans Mono}[Scale=0.88]
\setmonofont{DejaVu Sans Mono}[Scale=0.88]

% ---------- Page geometry ----------
\usepackage[
  letterpaper,
  top=1in,
  bottom=1in,
  left=1.25in,
  right=1.25in,
  headheight=14pt
]{geometry}

% ---------- Color + callout engine (must load BEFORE color-scheme input) ----
\usepackage{xcolor}
\usepackage[most]{tcolorbox}

% ---------- Typesetting packages ----------
\usepackage{tabularx}
\usepackage{longtable}
\usepackage{booktabs}
\usepackage{colortbl}
\usepackage{lastpage}
\usepackage{enumitem}
\usepackage{listings}
\usepackage{tikz}
\usepackage{fancyhdr}
\usepackage{titlesec}
\usepackage{parskip}
\usepackage{microtype}

% ---------- Pull in the color scheme (defines gsdnavy, gsdaccent, callouts) -
% ============================================================================
%  gsd-color-scheme.tex
%  Color palette + tcolorbox environment definitions for all GSD User Guides
% ============================================================================

\definecolor{gsdnavy}{HTML}{1B2A4A}
\definecolor{gsdaccent}{HTML}{CB3837}
\definecolor{gsdcodebg}{HTML}{F5F5F5}
\definecolor{gsdcodeborder}{HTML}{CCCCCC}

\definecolor{gsdtipbg}{HTML}{E8F4FD}
\definecolor{gsdtipbar}{HTML}{2B6CB0}

\definecolor{gsdwarnbg}{HTML}{FFF3CD}
\definecolor{gsdwarnbar}{HTML}{B78D00}

\definecolor{gsddangerbg}{HTML}{F8D7DA}
\definecolor{gsddangerbar}{HTML}{B02A37}

\definecolor{gsdsuccessbg}{HTML}{D4EDDA}
\definecolor{gsdsuccessbar}{HTML}{198754}

\definecolor{gsdnotebg}{HTML}{F0F0F5}
\definecolor{gsdnotebar}{HTML}{495057}

% ---------- Callout boxes ----------
\tcbset{
  gsdcalloutbase/.style={
    enhanced,
    breakable,
    left=10pt,
    right=8pt,
    top=6pt,
    bottom=6pt,
    boxrule=0pt,
    frame hidden,
    borderline west={4pt}{0pt}{#1},
    colback=#1,
    fonttitle=\bfseries\headingfont,
    fontupper=\small,
    arc=2pt
  }
}

\newtcolorbox{gsdtip}[1][Tip]{%
  gsdcalloutbase=gsdtipbg,
  borderline west={4pt}{0pt}{gsdtipbar},
  colback=gsdtipbg,
  title={#1}
}

\newtcolorbox{gsdwarning}[1][Warning]{%
  gsdcalloutbase=gsdwarnbg,
  borderline west={4pt}{0pt}{gsdwarnbar},
  colback=gsdwarnbg,
  title={#1}
}

\newtcolorbox{gsdimportant}[1][Important]{%
  gsdcalloutbase=gsddangerbg,
  borderline west={4pt}{0pt}{gsddangerbar},
  colback=gsddangerbg,
  title={#1}
}

\newtcolorbox{gsdnote}[1][Note]{%
  gsdcalloutbase=gsdnotebg,
  borderline west={4pt}{0pt}{gsdnotebar},
  colback=gsdnotebg,
  title={#1}
}

\newtcolorbox{gsdsuccess}[1][Success]{%
  gsdcalloutbase=gsdsuccessbg,
  borderline west={4pt}{0pt}{gsdsuccessbar},
  colback=gsdsuccessbg,
  title={#1}
}


% ---------- Hyperref (needs colors already defined) ----------
\usepackage[
  colorlinks=true,
  linkcolor=gsdnavy,
  urlcolor=gsdaccent,
  citecolor=gsdnavy,
  bookmarksnumbered=true
]{hyperref}

% ---------- Header / footer ----------
\pagestyle{fancy}
\fancyhf{}
\fancyhead[L]{\small\headingfont\@title}
\fancyhead[R]{\small\headingfont\leftmark}
\fancyfoot[L]{\small\headingfont GSD v1.35.0 — CC BY-SA 4.0}
\fancyfoot[R]{\small\headingfont Page \thepage\ of \pageref*{LastPage}}
\renewcommand{\headrulewidth}{0.4pt}
\renewcommand{\footrulewidth}{0.4pt}

\fancypagestyle{plain}{%
  \fancyhf{}
  \fancyfoot[L]{\small\headingfont GSD v1.35.0 — CC BY-SA 4.0}
  \fancyfoot[R]{\small\headingfont Page \thepage\ of \pageref*{LastPage}}
  \renewcommand{\headrulewidth}{0pt}
  \renewcommand{\footrulewidth}{0.4pt}
}

% ---------- Chapter / section styling ----------
\titleformat{\chapter}[display]
  {\headingfont\huge\bfseries\color{gsdnavy}}
  {\chaptertitlename\ \thechapter}{20pt}{\Huge}
\titleformat{\section}
  {\headingfont\Large\bfseries\color{gsdnavy}}
  {\thesection}{1em}{}
\titleformat{\subsection}
  {\headingfont\large\bfseries\color{gsdnavy}}
  {\thesubsection}{1em}{}

% ---------- Listings (code blocks) ----------
\lstdefinestyle{gsdcode}{
  basicstyle=\ttfamily\small,
  backgroundcolor=\color{gsdcodebg},
  frame=leftline,
  framerule=2pt,
  rulecolor=\color{gsdaccent},
  xleftmargin=1em,
  xrightmargin=0.5em,
  aboveskip=0.8em,
  belowskip=0.8em,
  breaklines=true,
  breakatwhitespace=true,
  showstringspaces=false,
  columns=flexible,
  keepspaces=true
}
\lstset{style=gsdcode}

% ---------- Semantic macros ----------
\newcommand{\cclicense}{%
  \textit{Released under CC BY-SA 4.0.\\
  Authored by Tibsfox. GSD itself is MIT-licensed — see gsd-build/get-shit-done.}%
}

\newcommand{\chapterheading}[1]{{\headingfont\LARGE\bfseries\color{gsdnavy}#1\par\medskip}}
\newcommand{\sectionheading}[1]{{\headingfont\Large\bfseries\color{gsdnavy}#1\par\smallskip}}
 in their
%  preamble (this file pulls in the color scheme transitively).
% ============================================================================

% ---------- Fonts ----------
\usepackage{fontspec}
\setmainfont{DejaVu Serif}[
  Ligatures=TeX,
  Scale=0.95,
  BoldFont={DejaVu Serif Bold},
  ItalicFont={DejaVu Serif Italic},
  BoldItalicFont={DejaVu Serif Bold Italic}
]
\newfontfamily\headingfont{DejaVu Sans}[Scale=0.95]
\newfontfamily\monofont{DejaVu Sans Mono}[Scale=0.88]
\setmonofont{DejaVu Sans Mono}[Scale=0.88]

% ---------- Page geometry ----------
\usepackage[
  letterpaper,
  top=1in,
  bottom=1in,
  left=1.25in,
  right=1.25in,
  headheight=14pt
]{geometry}

% ---------- Color + callout engine (must load BEFORE color-scheme input) ----
\usepackage{xcolor}
\usepackage[most]{tcolorbox}

% ---------- Typesetting packages ----------
\usepackage{tabularx}
\usepackage{longtable}
\usepackage{booktabs}
\usepackage{colortbl}
\usepackage{lastpage}
\usepackage{enumitem}
\usepackage{listings}
\usepackage{tikz}
\usepackage{fancyhdr}
\usepackage{titlesec}
\usepackage{parskip}
\usepackage{microtype}

% ---------- Pull in the color scheme (defines gsdnavy, gsdaccent, callouts) -
% ============================================================================
%  gsd-color-scheme.tex
%  Color palette + tcolorbox environment definitions for all GSD User Guides
% ============================================================================

\definecolor{gsdnavy}{HTML}{1B2A4A}
\definecolor{gsdaccent}{HTML}{CB3837}
\definecolor{gsdcodebg}{HTML}{F5F5F5}
\definecolor{gsdcodeborder}{HTML}{CCCCCC}

\definecolor{gsdtipbg}{HTML}{E8F4FD}
\definecolor{gsdtipbar}{HTML}{2B6CB0}

\definecolor{gsdwarnbg}{HTML}{FFF3CD}
\definecolor{gsdwarnbar}{HTML}{B78D00}

\definecolor{gsddangerbg}{HTML}{F8D7DA}
\definecolor{gsddangerbar}{HTML}{B02A37}

\definecolor{gsdsuccessbg}{HTML}{D4EDDA}
\definecolor{gsdsuccessbar}{HTML}{198754}

\definecolor{gsdnotebg}{HTML}{F0F0F5}
\definecolor{gsdnotebar}{HTML}{495057}

% ---------- Callout boxes ----------
\tcbset{
  gsdcalloutbase/.style={
    enhanced,
    breakable,
    left=10pt,
    right=8pt,
    top=6pt,
    bottom=6pt,
    boxrule=0pt,
    frame hidden,
    borderline west={4pt}{0pt}{#1},
    colback=#1,
    fonttitle=\bfseries\headingfont,
    fontupper=\small,
    arc=2pt
  }
}

\newtcolorbox{gsdtip}[1][Tip]{%
  gsdcalloutbase=gsdtipbg,
  borderline west={4pt}{0pt}{gsdtipbar},
  colback=gsdtipbg,
  title={#1}
}

\newtcolorbox{gsdwarning}[1][Warning]{%
  gsdcalloutbase=gsdwarnbg,
  borderline west={4pt}{0pt}{gsdwarnbar},
  colback=gsdwarnbg,
  title={#1}
}

\newtcolorbox{gsdimportant}[1][Important]{%
  gsdcalloutbase=gsddangerbg,
  borderline west={4pt}{0pt}{gsddangerbar},
  colback=gsddangerbg,
  title={#1}
}

\newtcolorbox{gsdnote}[1][Note]{%
  gsdcalloutbase=gsdnotebg,
  borderline west={4pt}{0pt}{gsdnotebar},
  colback=gsdnotebg,
  title={#1}
}

\newtcolorbox{gsdsuccess}[1][Success]{%
  gsdcalloutbase=gsdsuccessbg,
  borderline west={4pt}{0pt}{gsdsuccessbar},
  colback=gsdsuccessbg,
  title={#1}
}


% ---------- Hyperref (needs colors already defined) ----------
\usepackage[
  colorlinks=true,
  linkcolor=gsdnavy,
  urlcolor=gsdaccent,
  citecolor=gsdnavy,
  bookmarksnumbered=true
]{hyperref}

% ---------- Header / footer ----------
\pagestyle{fancy}
\fancyhf{}
\fancyhead[L]{\small\headingfont\@title}
\fancyhead[R]{\small\headingfont\leftmark}
\fancyfoot[L]{\small\headingfont GSD v1.35.0 — CC BY-SA 4.0}
\fancyfoot[R]{\small\headingfont Page \thepage\ of \pageref*{LastPage}}
\renewcommand{\headrulewidth}{0.4pt}
\renewcommand{\footrulewidth}{0.4pt}

\fancypagestyle{plain}{%
  \fancyhf{}
  \fancyfoot[L]{\small\headingfont GSD v1.35.0 — CC BY-SA 4.0}
  \fancyfoot[R]{\small\headingfont Page \thepage\ of \pageref*{LastPage}}
  \renewcommand{\headrulewidth}{0pt}
  \renewcommand{\footrulewidth}{0.4pt}
}

% ---------- Chapter / section styling ----------
\titleformat{\chapter}[display]
  {\headingfont\huge\bfseries\color{gsdnavy}}
  {\chaptertitlename\ \thechapter}{20pt}{\Huge}
\titleformat{\section}
  {\headingfont\Large\bfseries\color{gsdnavy}}
  {\thesection}{1em}{}
\titleformat{\subsection}
  {\headingfont\large\bfseries\color{gsdnavy}}
  {\thesubsection}{1em}{}

% ---------- Listings (code blocks) ----------
\lstdefinestyle{gsdcode}{
  basicstyle=\ttfamily\small,
  backgroundcolor=\color{gsdcodebg},
  frame=leftline,
  framerule=2pt,
  rulecolor=\color{gsdaccent},
  xleftmargin=1em,
  xrightmargin=0.5em,
  aboveskip=0.8em,
  belowskip=0.8em,
  breaklines=true,
  breakatwhitespace=true,
  showstringspaces=false,
  columns=flexible,
  keepspaces=true
}
\lstset{style=gsdcode}

% ---------- Semantic macros ----------
\newcommand{\cclicense}{%
  \textit{Released under CC BY-SA 4.0.\\
  Authored by Tibsfox. GSD itself is MIT-licensed — see gsd-build/get-shit-done.}%
}

\newcommand{\chapterheading}[1]{{\headingfont\LARGE\bfseries\color{gsdnavy}#1\par\medskip}}
\newcommand{\sectionheading}[1]{{\headingfont\Large\bfseries\color{gsdnavy}#1\par\smallskip}}


% ============================================================================
%  Going-Further-specific preamble extensions
% ============================================================================

% ---------- Part page styling ----------
\titleformat{\part}[display]
  {\headingfont\Huge\bfseries\color{gsdnavy}}
  {\partname\ \thepart}{20pt}{\Huge}

% ---------- Scenario walkthrough environment (the signature element) ------
\definecolor{gsdscenariobg}{HTML}{F9F9F9}

\newtcolorbox{gsdscenario}[1][Scenario]{%
  enhanced,
  breakable,
  left=12pt,
  right=10pt,
  top=8pt,
  bottom=8pt,
  boxrule=0pt,
  frame hidden,
  borderline west={6pt}{0pt}{gsdaccent},
  colback=gsdscenariobg,
  fonttitle=\bfseries\itshape\headingfont,
  fontupper=\small,
  arc=2pt,
  title={Scenario: #1}
}

% ---------- Per-language code listing presets -----------------------------
\lstdefinelanguage{jsonl}{
  literate=
    *{0}{{{\color{gsdnavy}0}}}{1}
     {1}{{{\color{gsdnavy}1}}}{1}
     {2}{{{\color{gsdnavy}2}}}{1}
     {3}{{{\color{gsdnavy}3}}}{1}
     {4}{{{\color{gsdnavy}4}}}{1}
     {5}{{{\color{gsdnavy}5}}}{1}
     {6}{{{\color{gsdnavy}6}}}{1}
     {7}{{{\color{gsdnavy}7}}}{1}
     {8}{{{\color{gsdnavy}8}}}{1}
     {9}{{{\color{gsdnavy}9}}}{1}
     {:}{{{\color{gsdaccent}{:}}}}{1}
     {,}{{{\color{gsdaccent}{,}}}}{1}
     {\{}{{{\color{gsdaccent}{\{}}}}{1}
     {\}}{{{\color{gsdaccent}{\}}}}}{1}
     {[}{{{\color{gsdaccent}{[}}}}{1}
     {]}{{{\color{gsdaccent}{]}}}}{1},
  string=[b]",
  stringstyle=\color{gsdnavy!80!black},
  showstringspaces=false
}

\lstdefinestyle{gsdjsonstyle}{
  style=gsdcode,
  language=jsonl,
}

\lstnewenvironment{gsdjson}{\lstset{style=gsdjsonstyle}}{}

\lstdefinestyle{gsdbashstyle}{
  style=gsdcode,
  language=bash,
  morekeywords={npx,claude,gsd,git,cd,export,source,echo,cat,grep,ls,mkdir,rm,touch,curl,wget},
  keywordstyle=\color{gsdnavy}\bfseries,
  commentstyle=\color{gray}\itshape,
}

\lstnewenvironment{gsdbash}{\lstset{style=gsdbashstyle}}{}

\lstdefinelanguage{yamlflavor}{
  morekeywords={true,false,null,yes,no},
  keywordstyle=\color{gsdaccent},
  morecomment=[l]{\#},
  commentstyle=\color{gray}\itshape,
  morestring=[b]',
  morestring=[b]",
  stringstyle=\color{gsdnavy!80!black},
}

\lstdefinestyle{gsdyamlstyle}{
  style=gsdcode,
  language=yamlflavor,
}

\lstnewenvironment{gsdyaml}{\lstset{style=gsdyamlstyle}}{}

\lstnewenvironment{gsdplain}{\lstset{style=gsdcode,language=}}{}

% ---------- Reusable credential safety callout ----------------------------
\newcommand{\credentialwarning}{%
\begin{gsdwarning}[Credential Safety]
The values shown above are \textbf{placeholders}. Never commit real API keys,
OAuth tokens, or database connection strings to a file that will be checked
into version control. Use environment variables (\texttt{\$\{VAR\_NAME\}}) and
add \texttt{.mcp.json} to \texttt{.gitignore} whenever it carries credentials.
\end{gsdwarning}%
}

\begin{document}

% ---------- Title page ----------
\begin{titlepage}
\centering
\vspace*{2cm}
{\headingfont\Huge\bfseries\color{gsdnavy} GSD\par}
\vspace{0.4cm}
{\headingfont\LARGE\bfseries\color{gsdnavy} Going Further\par}
\vspace{0.8cm}
{\headingfont\large Skills, MCP, and Multi-Repo Workflows\par}
\vspace{2cm}
\begin{tcolorbox}[
  enhanced,
  width=0.78\textwidth,
  colback=gsdnotebg,
  colframe=gsdnotebar,
  boxrule=0.6pt,
  arc=3pt,
  fonttitle=\bfseries\headingfont,
  title={What This Guide Covers}
]
\small
Three advanced GSD topics, each delivered through a complete walkthrough you
can follow at the keyboard:
\begin{itemize}[leftmargin=1.2em,itemsep=0.15em]
  \item \textbf{Skills \& Slash Commands} --- the SKILL.md format, the
        12-runtime install matrix, and how to inject custom domain
        knowledge into GSD's subagents.
  \item \textbf{MCP Integration} --- connecting GSD to external tools via
        the Model Context Protocol, with full walkthroughs for Google
        Stitch UI generation and schema-aware planning via a Postgres MCP.
  \item \textbf{Multi-Repo Workflows} --- workstreams, workspaces, the
        STATE.md ownership rule, and patterns for teams running GSD across
        microservice fleets.
\end{itemize}
\end{tcolorbox}
\vspace{2cm}
{\large Tibsfox\par}
\vspace{0.4cm}
{\small For GSD v1.35.0 \par}
\vspace{0.4cm}
{\small Audience: GSD users ready to customize, integrate, and scale \par}
\vfill
{\footnotesize \cclicense\par}
\end{titlepage}

\tableofcontents
\clearpage

% ============================================================================
%  part1-skills.tex
%  Part I of "GSD: Going Further" — Skills & Slash Commands
%  Input into the master document via \input{part1-skills}
% ============================================================================

\part{Skills \& Slash Commands}

\chapter{How GSD Actually Works --- Skills Anatomy}

Most GSD users never look under the hood. They install the tool, type
\texttt{/gsd:new-project}, and watch Claude Code do the rest. That is exactly
how it is supposed to feel. But the moment you want to bend GSD to fit your
team, your domain, or your personal taste, the black box stops being a
feature and starts being an obstacle.

This chapter opens the box. By the end of it you will know where GSD stores
its skills on disk, what a single skill file looks like, how Claude Code
decides which skills to activate, and why the installer ships two different
layouts depending on which version of Claude Code you have. None of this
requires source code spelunking --- everything GSD installs is plain
markdown, which means you can read it, diff it, copy it, and edit it with
the same tools you already use.

\section{Finding the Skills on Disk}

On a Claude Code 2.1.88 or newer install, GSD writes its skills into your
per-user Claude configuration directory. On Linux and macOS that lives at
\texttt{\~{}/.claude/skills/}; on Windows it is
\texttt{\%USERPROFILE\%\textbackslash.claude\textbackslash skills\textbackslash}.
The GSD installer creates one sub-directory per command, each prefixed with
\texttt{gsd-}, and each containing a single \texttt{SKILL.md} file:

\begin{gsdplain}
~/.claude/skills/
├── gsd-new-project/
│   └── SKILL.md
├── gsd-plan-phase/
│   └── SKILL.md
├── gsd-execute-phase/
│   └── SKILL.md
├── gsd-verify-work/
│   └── SKILL.md
├── gsd-progress/
│   └── SKILL.md
├── gsd-help/
│   └── SKILL.md
└── ... (30+ more)
\end{gsdplain}

Go ahead and navigate there now. Run \texttt{ls \~{}/.claude/skills/} on
your own machine. You should see a long list of \texttt{gsd-*} directories
plus any other skills Claude Code has picked up from other tools or from
your own hand-authored experiments. This is not a sandbox --- it is your
real home directory, and those markdown files are the exact instructions
Claude reads when you invoke a GSD command.

\begin{gsdtip}
You can read any GSD skill file. They are plain markdown --- no compilation,
no opaque binaries, no generated code. If you are ever unsure what a
\texttt{/gsd:*} command is \emph{actually} going to do, open its
\texttt{SKILL.md} in your editor and read it. That is the entire contract.
\end{gsdtip}

\section{The SKILL.md Format}

Every skill file Claude Code 2.1.88 and later consumes follows the same
shape: a YAML frontmatter block, a heading, and an instruction body. The
frontmatter carries two required keys --- \texttt{name} and
\texttt{description} --- and the body is markdown prose that becomes the
instruction set the model sees when the skill activates.

Here is the minimum viable skill file:

\begin{gsdplain}
---
name: gsd-progress
description: "Show GSD project progress and current phase status"
---
# GSD Progress Report

When the user asks about the current state of the project, read
`.planning/STATE.md` and `.planning/ROADMAP.md`. Produce a compact
summary containing:

1. The current milestone (from ROADMAP.md).
2. The active phase (from STATE.md).
3. The most recent three completed work items.
4. Any open blockers recorded under the "Blockers" heading.

Render the output as a markdown table with two columns: "Item" and
"Status". Do not fabricate phase data --- if STATE.md is missing, say
so and suggest running `/gsd:new-project` to initialize it.
\end{gsdplain}

That is a complete, working skill. Drop that file into
\texttt{\~{}/.claude/skills/gsd-progress/SKILL.md}, reopen Claude Code, and
typing \textit{``show me the GSD progress''} will surface it as a candidate
skill. The \texttt{name} field is the stable identifier --- it also
determines the slash-command alias (\texttt{/gsd:progress} maps to
\texttt{gsd-progress}) on runtimes that expose slash commands. The
\texttt{description} field is where the real magic lives, and it deserves
its own section.

\section{How Trigger Descriptions Work}

Claude Code does not load every skill on disk into every request. That
would blow the context window inside of a dozen commands. Instead the
runtime keeps a lightweight manifest of all installed skills --- just their
names and descriptions --- and when a new user message arrives, it pattern
matches the message against those descriptions to decide which skills to
promote into the full context. Only promoted skills have their
instruction body loaded.

The \texttt{description} is therefore the \emph{activation contract}. A
vague description like \texttt{"Help with GSD"} will match almost
everything and almost nothing usefully. A concrete one like
\texttt{"Show GSD project progress and current phase status"} will match
phrases like \textit{``where are we on the project''},
\textit{``what phase are we in''}, or \textit{``give me a status update''}
without swamping unrelated requests.

When you author your own GSD-adjacent skills, invest time in the
description line. Use action verbs, name the artifact you operate on, and
include a few of the phrasings users actually say out loud. A good
description is not a summary of what the skill does --- it is a list of
\emph{triggers} that should fire it.

\begin{gsdnote}
The frontmatter is strict YAML. Unquoted colons in the description will
break the parser and silently disable the skill. When in doubt, wrap the
description in double quotes. This is the single most common reason a
hand-authored skill ``doesn't work''.
\end{gsdnote}

\section{The Installer's Role}

GSD ships as an npm package. When you run \texttt{npx gsd-build install},
the installer does three things in sequence:

\begin{enumerate}
\item \textbf{Detect the runtime.} It looks at which agent CLIs are
installed on the current machine --- Claude Code, Codex, Cursor, Windsurf,
and so on --- and also at the \emph{version} of Claude Code it finds. The
2.1.88 cutoff matters: earlier versions do not understand the skills
directory layout at all.
\item \textbf{Read its source skill definitions.} GSD carries its command
definitions in a single canonical form inside the npm package. These are
the master copies.
\item \textbf{Transform and write.} For each detected runtime, the
installer rewrites the master skill into whatever format that runtime
actually consumes, then drops the result into the right directory. One
master, many derivatives.
\end{enumerate}

This architecture is why the \emph{same} \texttt{/gsd:plan-phase} command
behaves identically whether you are driving Claude Code, Codex, or Cursor
--- the installer is doing the translation work up front so the user-facing
surface is uniform.

\section{Skills versus Legacy Slash Commands}

Before Claude Code 2.1.88, Claude Code had no notion of skills. The only
way to ship reusable instructions was as a \emph{slash command}: a
markdown file living at \texttt{\~{}/.claude/commands/gsd/[name].md} that
the user invoked by literally typing its name. There was no description
field, no trigger matching, no auto-activation --- the user had to know
the command existed and type it explicitly.

GSD still supports this old world. If the installer detects a pre-2.1.88
Claude Code install, it writes legacy slash commands instead of skills:

\begin{gsdplain}
~/.claude/commands/gsd/
├── new-project.md
├── plan-phase.md
├── execute-phase.md
├── verify-work.md
└── ...
\end{gsdplain}

You interact with both layouts the same way from the user side ---
\texttt{/gsd:plan-phase} works whether it resolves to a skill or a legacy
command --- but under the hood the runtime treats them differently. Skills
are auto-activating and context-aware; legacy commands only fire on an
explicit invocation. If you upgrade Claude Code from an older version, run
\texttt{npx gsd-build install} again and the installer will migrate you
forward automatically, leaving the legacy directory in place so nothing
breaks mid-session.

\section{What You Can Do With This Knowledge}

Once you internalize that GSD skills are just markdown on disk, a long
list of previously scary tasks becomes trivial. You can diff two installs
to see what changed between GSD releases. You can keep a local patch
directory and re-apply your own edits after an upgrade. You can fork a
skill, rename it to \texttt{gsd-plan-phase-health}, tailor it to your
domain, and have both the generic and the specialized version coexist in
the same directory. You can hand-write skills that call GSD primitives
without waiting for upstream to ship them. Everything is text, everything
is readable, and nothing is hidden.

The rest of this part walks through each of those capabilities in detail,
starting with the compatibility matrix that determines where your skills
actually land.

% ============================================================================
\chapter{Skills Across Runtimes}

GSD is not a Claude Code exclusive. The installer knows about a dozen
different agent runtimes, each with its own directory layout, its own
file format, and its own opinions about how instructions should be
delivered to the model. This chapter is the reference you will reach for
when a GSD command that works on one machine appears to be missing on
another. The short answer is almost always \textit{``it installed fine,
but it lives somewhere else on this runtime.''}

\section{The Compatibility Matrix}

The table below lists every runtime the GSD installer currently supports,
the format it emits for that runtime, the directory where the output
lands, and the command you can run to confirm the install succeeded.

{\footnotesize
\begin{longtable}{@{}p{2.5cm} p{3.2cm} p{4.2cm} p{2.8cm}@{}}
\toprule
\textbf{Runtime} & \textbf{Install Format} & \textbf{Directory (Global)} & \textbf{Verify Command} \\
\midrule
\endfirsthead
\toprule
\textbf{Runtime} & \textbf{Install Format} & \textbf{Directory (Global)} & \textbf{Verify Command} \\
\midrule
\endhead
Claude Code 2.1.88+ & \texttt{skills/gsd-*/SKILL.md} & \texttt{\~{}/.claude/skills/} & \texttt{/gsd:help} \\
Claude Code (legacy) & \texttt{commands/gsd/*.md} & \texttt{\~{}/.claude/commands/} & \texttt{/gsd:help} \\
Codex & \texttt{skills/gsd-*/SKILL.md} & \texttt{\~{}/.codex/skills/} & \texttt{\$gsd-help} \\
Copilot & prompts + agents & \texttt{\~{}/.github/} & \texttt{/gsd:help} \\
Cursor & \texttt{skills/gsd-*/SKILL.md} & \texttt{\~{}/.cursor/} & \texttt{/gsd:help} \\
Windsurf & markdown transform & \texttt{\~{}/.windsurf/} & \texttt{/gsd:help} \\
OpenCode & config file & \texttt{\~{}/.config/opencode/} & \texttt{/gsd-help} \\
Gemini CLI & skills & \texttt{\~{}/.gemini/} & \texttt{/gsd:help} \\
Antigravity & skills & \texttt{\~{}/.gemini/antigravity/} or \texttt{./.agent/} & \texttt{/gsd:help} \\
Cline & \texttt{.clinerules} & \texttt{\~{}/.cline/} & rules auto-loaded \\
Augment & skills transform & \texttt{\~{}/.augment/} & \texttt{/gsd:help} \\
Qwen Code & skills (open standard) & \texttt{\~{}/.qwen/skills/} & \texttt{/gsd:help} \\
\bottomrule
\end{longtable}
}

Twelve runtimes, twelve directory layouts. Four of them (Claude Code 2.1.88
and later, Codex, Cursor, and Qwen Code) converge on the same skills-with-
SKILL.md shape, which is the format this book treats as canonical. The
other eight each require the installer to do some amount of translation
work, and understanding those translations is the key to debugging any
cross-runtime weirdness you run into.

\section{How the Installer Transforms}

The installer maintains a single internal representation for each GSD
command --- call it the ``source form.'' When a runtime requires a
different format, the installer applies a translation pass. Here are the
three most important translations.

\subsection{Skills into Prompts (Copilot)}

GitHub Copilot does not have a skills directory. It has a
\texttt{prompts/} directory for standalone prompt files and an
\texttt{agents/} directory for role definitions, and the two work
together. The installer splits each GSD skill in half: the instruction
body becomes a prompt file under \texttt{\~{}/.github/prompts/gsd/}, and
the trigger metadata plus any agent role information becomes an agent
definition under \texttt{\~{}/.github/agents/}. Invoking
\texttt{/gsd:plan-phase} in Copilot fires both sides simultaneously.

\subsection{Skills into Rules (Cline)}

Cline takes the opposite approach. It has no concept of slash commands at
all --- instead, it reads a \texttt{.clinerules} file and treats the
entire contents as a persistent instruction prefix injected into every
request. The installer concatenates every GSD skill body into a single
\texttt{.clinerules} file and drops it at
\texttt{\~{}/.cline/.clinerules}. You will not find a
\texttt{/gsd:plan-phase} slash command in Cline because there are no
slash commands to find; the rules simply load automatically and Cline
recognizes the GSD phrasings when you use them conversationally.

\subsection{Skills into TOML + Markdown (Codex)}

Codex uses a paired-file convention: a \texttt{.toml} file declares the
skill metadata, and a matching \texttt{.md} file alongside it carries the
instruction body. The installer emits both halves. If you open
\texttt{\~{}/.codex/skills/gsd-plan-phase/} you will see a
\texttt{skill.toml} with the name and description fields, plus a
\texttt{SKILL.md} with the prose. Codex also uses a different invocation
prefix: \texttt{\$gsd-help} instead of \texttt{/gsd:help}, which is worth
memorizing if you move between machines.

\section{Runtime-Specific Gotchas}

Even when you know the translation rules, a few runtimes have quirks
worth calling out explicitly.

\subsection{Cline Has No Slash Commands}

If a new Cline user complains that \texttt{/gsd:plan-phase} ``isn't
working,'' they have usually not done anything wrong --- Cline literally
does not support slash commands. The fix is not to reinstall. The fix is
to describe what you want in natural language (\textit{``plan the next
phase''}) and let the injected \texttt{.clinerules} do its job.

\subsection{Copilot Needs Both Halves}

The prompt and agent files Copilot writes work together, and if only one
half lands (for example, because the installer was interrupted or because
\texttt{\~{}/.github/agents/} already contained a conflicting file), the
GSD command will misbehave in subtle ways. Re-running
\texttt{npx gsd-build install --force} will overwrite any half-written
state and is the standard fix.

\subsection{Codex Uses TOML Pairs}

If you hand-edit a Codex skill, remember that the metadata lives in the
\texttt{.toml} file, not in frontmatter. Editing the \texttt{SKILL.md}
body to change the \texttt{description} field will have no effect ---
Codex never reads frontmatter. Edit the \texttt{.toml}.

\subsection{Antigravity Has Two Install Dirs}

Antigravity supports both a global install at
\texttt{\~{}/.gemini/antigravity/} and a project-local install at
\texttt{./.agent/}. The installer picks global by default but will honor
the \texttt{--local} flag to install into the current project. This is
the only runtime where GSD supports per-project scoped command installs
at the runtime layer; for every other runtime, per-project customization
happens through the \texttt{agent\_skills} mechanism described in the
next chapter.

\section{Non-Claude Model Profiles}

Most GSD subagents (the executor, planner, researcher, and verifier
roles) are built on Claude Code's subagent spawn primitive, which
defaults to whichever Anthropic model the parent session is running.
That is usually what you want --- consistency across the session --- but
it is a footgun the first time you run GSD inside an OpenRouter-fronted
Claude Code or on a locally hosted model.

The problem: even when the parent session is routed through OpenRouter,
a spawned subagent may try to call Anthropic directly unless you
explicitly tell it not to. You end up paying OpenRouter for the parent
and Anthropic for every subagent spawn simultaneously.

GSD ships an \texttt{inherit} profile for exactly this case. When you
select the \texttt{inherit} profile (via \texttt{/gsd:set-profile
inherit} or by editing \texttt{.planning/config.json}), every spawned
subagent will reuse the parent session's model configuration instead of
defaulting to Anthropic. If the parent is on OpenRouter, the subagents
are on OpenRouter. If the parent is on a local Ollama model, the
subagents are too.

\begin{gsdimportant}
If GSD subagents call Anthropic models and you are paying through
OpenRouter, switch to the \texttt{inherit} profile or you will see
double-billing. Run \texttt{/gsd:set-profile inherit} once per project.
\end{gsdimportant}

\section{Verifying an Install}

Regardless of runtime, the fastest way to confirm GSD is actually
installed is to run the \texttt{/gsd:help} command (or its runtime-specific
equivalent from the matrix). It is a skill like any other, and if it
responds with the GSD command listing, everything else is wired up
correctly. If it does not respond, the first diagnostic step is to
\texttt{ls} the install directory for your runtime and confirm the
\texttt{gsd-*} entries exist. Ninety percent of install problems are
really ``I am on a runtime I did not expect to be on'' problems, and the
matrix is there to cut that debugging time down to seconds.

% ============================================================================
\chapter{Agent Skills Injection}

So far everything in this Part has been about \emph{global} skills --- the
files GSD installs into your per-user directory that apply to every
project you work on. That is the right default for the command layer.
But GSD also has a sharper tool for per-project customization, one that
does not require touching the global directory at all: \texttt{agent\_skills}.

\texttt{agent\_skills} is the mechanism by which a single project teaches
its GSD subagents things about its own domain. Those teachings live
inside the project, travel with it through version control, and apply
only while you are working in that project. No other project sees them.
No global state gets polluted.

\section{The Concept}

When GSD spawns a subagent --- an executor running a phase, a planner
drafting the next set of work items, or a researcher collecting
domain information --- it builds the subagent's system prompt from a
stack of pieces: the subagent's role definition, the current project
state, the task at hand, and optionally, one or more
\emph{injected skill blocks} drawn from local project directories.

Those injected blocks arrive in the subagent's prompt as an
\texttt{<agent\_skills>} section. From the subagent's point of view,
they are just more instructions --- indistinguishable from the built-in
role definition except that they came from a project-local directory
rather than from the GSD source tree. This is pure prompt augmentation.
There is no plugin system, no sandbox, no new execution path. The
injected content simply becomes part of what the subagent sees.

That simplicity is the feature. You can teach a GSD researcher about
FDA endpoints the same way you would teach a human researcher about
them: by handing over a short written brief. The researcher reads the
brief, integrates it into its working knowledge, and uses it for the
duration of the task.

\section{Configuration Syntax}

\texttt{agent\_skills} is configured in the project's
\texttt{.planning/config.json} file. The key is \texttt{agent\_skills},
and its value is an object mapping subagent type names to arrays of
directory paths. Each path is resolved relative to the project root.

\begin{gsdjson}
{
  "agent_skills": {
    "executor": ["./skills/domain-knowledge/"],
    "planner": ["./skills/planning-rules/"],
    "researcher": ["./skills/research-sources/"]
  }
}
\end{gsdjson}

When a subagent of a given type spawns, the installer walks the
directories listed under that type, reads every \texttt{.md} file it
finds, concatenates the contents, and wraps the result in an
\texttt{<agent\_skills>} block inside the subagent's prompt. Files are
read in lexicographic order, so prefixing names with \texttt{01-},
\texttt{02-}, and so on gives you explicit control over the final
ordering when order matters.

\section{Which Agent Types Accept Skills}

Not every subagent type honors \texttt{agent\_skills}. As of the current
release the supported types are:

\begin{itemize}
\item \textbf{executor} --- the phase-running worker. Injected skills
here shape how code gets written, which conventions get followed, and
which libraries get preferred. Good for ``our codebase uses X pattern''
style rules.
\item \textbf{planner} --- the task decomposition agent. Injected skills
here shape how phases get broken into work items and what kinds of
safety rails the planner enforces. Good for ``never do X without Y''
style rules.
\item \textbf{researcher} --- the information-gathering agent. Injected
skills here shape which sources get consulted and what domain
vocabulary gets used in the resulting \texttt{RESEARCH.md}. Good for
teaching the researcher about APIs, standards, databases, and
terminology it would not otherwise know.
\end{itemize}

The verifier, documenter, and other specialized subagent types currently
use a fixed instruction set and do not read \texttt{agent\_skills}. That
may change in a future release; for now, if you need to customize
verification behavior, you do it by editing the global skill file
directly.

\section{Writing a Skill Directory}

A skill directory is just a directory. Create it anywhere under your
project root (convention is \texttt{./skills/[name]/}), drop one or
more \texttt{.md} files in it, and list the directory in
\texttt{agent\_skills}. That is the whole contract.

The markdown files inside do not need frontmatter. They do not need a
\texttt{name} or \texttt{description} field. They just need to contain
the instructions you want the subagent to read. You can write them in
whatever style feels natural --- prose, bullet lists, labeled sections,
annotated code examples --- and the subagent will parse them the same
way it parses any other instruction block.

The reason naming is descriptive (rather than formal) is that every
\texttt{.md} file's content gets concatenated together at spawn time.
The subagent does not see file boundaries; it sees one long block of
text. File names serve only two purposes: ordering (lexicographic) and
readability for you when you come back to the directory in six months.

\section{How Injection Affects Behavior}

The injected \texttt{<agent\_skills>} block lives inside the subagent's
prompt for the entire duration of its task. That means:

\begin{itemize}
\item The instructions persist across every tool call the subagent makes.
A researcher that spawns a web search, then a file read, then another
web search still has the FDA brief visible through all three calls.
\item The instructions are visible to the model at every decision point.
When the model is picking which URL to fetch or which section to write
first, it is weighing those choices against the injected skill content.
\item The instructions do not persist beyond the task. Once the subagent
exits, the injected block goes with it. The next subagent spawn rebuilds
the prompt from scratch, picking up any edits you made to the skill
files in between.
\end{itemize}

That last point is important: you can edit a skill file mid-session and
the next subagent spawn picks up the change immediately. There is no
rebuild step, no cache to invalidate, no installer to re-run.

\section{A Minimal Example}

Suppose you are building a product that must never merge to
\texttt{main} without first running an integration test. You want the
GSD planner to refuse to generate a plan that skips this step. Here is
the full contents of \texttt{./skills/planning-rules/no-untested-merges.md}:

\begin{gsdplain}
# Merge Safety Rule

Never produce a plan whose final step is a merge to `main`, `master`,
or any protected branch unless the plan also contains an explicit
integration test step that runs BEFORE the merge and must pass.

If a user asks for a plan that would merge without testing, respond
with a plan that adds the missing test step and a one-line note
explaining why.

Applies to: all merge-containing phases, all branches matching
`main`, `master`, `release/*`, or `prod`.
\end{gsdplain}

Six lines of instruction. That is enough. With this file in place and
\texttt{planner} pointing at \texttt{./skills/planning-rules/} in
\texttt{agent\_skills}, every subsequent \texttt{/gsd:plan-phase}
invocation in this project will generate plans that include a test
step before any merge, and will add a short justification when the
user's request would have skipped one.

The planner did not need to be recompiled. The GSD installer did not
need to be rerun. You wrote six lines of markdown, saved the file, and
the next planner spawn behaved differently. That is the entire feature.

\begin{gsdnote}
\texttt{agent\_skills} are project-local. They are defined in
\texttt{.planning/config.json}, they reference paths under the project
root, and they do not affect any other project on the same machine.
Commit the skill directories along with your code so the next developer
who clones the repo gets the same behavior automatically.
\end{gsdnote}

\section{When to Reach for This Tool}

\texttt{agent\_skills} is the right tool whenever your customization
need has three properties: it is specific to one project, it is worth
writing down, and it would benefit from being seen by a particular
subagent role on every spawn. Domain vocabulary, compliance rules,
in-house API conventions, preferred libraries, house coding standards,
data-handling constraints --- all of these fit naturally.

It is the wrong tool for global rules (those belong in a hand-authored
global skill), for one-off asks (those belong in the prompt itself),
and for rules the whole team has not yet agreed on (those belong in a
discussion, not a config file). Within its sweet spot, though, it is
the single most powerful lever GSD gives you for shaping subagent
behavior --- and the next chapter walks through a complete,
end-to-end use of it.

% ============================================================================
\chapter{Scenario --- Custom Research Skill}

\begin{gsdscenario}[Teaching GSD's Researchers Your Domain]

You are three weeks into building a small health-tech application --- a
patient intake form that needs to map its fields cleanly to the FHIR
Patient resource, check drug interactions against an authoritative
source, and keep one eye on the 510(k) clearance pathway in case the
product ever grows into a regulated medical device. You are using GSD
to manage the work. Planning, execution, and verification all run
smoothly. But the research phase keeps coming up short.

Here is what you want to fix, and how \texttt{agent\_skills} fixes it.

\textbf{\large The problem.}

You ran \texttt{/gsd:research-phase} yesterday on the topic
\textit{``data sources for drug interaction checking.''} The researcher
subagent did its job: it spawned, ran a batch of web searches, read a
handful of Wikipedia pages and generic developer blogs, and wrote a
tidy \texttt{RESEARCH.md} file into \texttt{.planning/research/}.
Tidy, but wrong for your use case.

The output was full of generic phrases like \textit{``drug interaction
APIs such as those offered by various healthcare data vendors,''}
hand-waving references to \textit{``the standard medical data format,''}
and not a single mention of the FDA's openFDA Drug Event API, the RxNorm
concept hierarchy, or the DrugBank identifier system. The section on
data formats talked about JSON generically but never named FHIR, never
mentioned the \texttt{MedicationStatement} resource, and never hinted
that a serious health-tech product would need to handle HL7 interop.

This is the researcher working exactly as designed. The problem is not
that it is bad at research; the problem is that it has no reason to
prefer FDA over DrugBank over a random consumer health blog, because
nobody told it which sources are authoritative in your domain. You are
going to tell it now.

\textbf{\large Step 1: Create the skill directory.}

Start by carving out a home for the domain brief. GSD convention is to
put project-local skill directories under \texttt{./skills/}, so make
one there with a descriptive name:

\begin{gsdbash}
mkdir -p ./skills/health-research/
touch ./skills/health-research/SKILL.md
\end{gsdbash}

Open that new \texttt{SKILL.md} in your editor and write a real brief.
This is the file you will hand to every future researcher spawn, so
make it concrete. Name the endpoints, name the resources, name the
regulations. Here is the full contents of the file we used for this
project:

\begin{gsdplain}
# Health-Tech Research Brief

When researching any topic related to drug data, patient records,
clinical workflows, or medical devices, prefer the following
authoritative sources over generic web results. Always cite them
by name in the final RESEARCH.md.

## Drug data
- FDA openFDA Drug Event API:
  https://api.fda.gov/drug/event.json
  Use for adverse event lookups and label data.
- RxNorm (NLM):
  https://rxnav.nlm.nih.gov/REST/
  Use for normalized drug name resolution.
- DrugBank (academic tier):
  https://go.drugbank.com/
  Use for interaction data; cite the DrugBank ID.

## Clinical trials and evidence
- ClinicalTrials.gov search API:
  https://clinicaltrials.gov/api/
  Use for active and completed trials, inclusion criteria,
  and results summaries.

## Data model
- FHIR R4 resource definitions:
  https://hl7.org/fhir/R4/resourcelist.html
  When modeling patient data, reference the exact FHIR resource
  and field names (e.g., Patient.identifier, Observation.code,
  MedicationStatement.medicationCodeableConcept).

## Regulatory
- HIPAA Privacy Rule: treat all PHI fields as requiring encryption
  at rest and in transit. Never store unmasked PHI in logs or
  telemetry.
- 510(k) Premarket Notification: if a feature could classify as
  a medical device, flag it and link to
  https://www.fda.gov/medical-devices/premarket-submissions/
  premarket-notification-510k
\end{gsdplain}

Notice what this file is \emph{not}. It is not a summary of everything
there is to know about health tech. It is not a tutorial. It is not a
policy document. It is a short, dense list of \textit{``here are the
sources to use and the vocabulary to use when writing up your
findings.''} Twenty-some lines is plenty. The researcher is already
capable of doing research; you are just handing it a better rolodex.

\textbf{\large Step 2: Configure \texttt{agent\_skills}.}

Open \texttt{.planning/config.json} and add a \texttt{researcher} entry
to the \texttt{agent\_skills} object. If the object does not yet exist,
create it. Here is the relevant snippet with the new entry in place:

\begin{gsdjson}
{
  "project_name": "health-intake",
  "profile": "inherit",
  "agent_skills": {
    "researcher": ["./skills/health-research/"]
  }
}
\end{gsdjson}

Save the file. That is the entire configuration step. No installer
rerun, no restart, no validation pass. The GSD runtime reads
\texttt{.planning/config.json} on every subagent spawn, so the next
time a researcher starts, it will pick up the new skill directory.

\textbf{\large Step 3: Run \texttt{/gsd:research-phase}.}

Re-run the same research phase you ran yesterday. This time the
planner subagent reads your request, kicks off a researcher spawn, and
the GSD runtime assembles the researcher's prompt: role definition,
task description, project state, and now --- because you wired it up
--- an \texttt{<agent\_skills>} block containing your health-tech brief.

\begin{gsdbash}
/gsd:research-phase "data sources for drug interaction checking"
\end{gsdbash}

What happens inside the spawn: the researcher reads the brief as part
of its initial context, sees the named endpoints (openFDA, RxNorm,
DrugBank), and queues those searches first. When it hits web results
for generic health blogs, it now has a yardstick to measure them
against, and it deprioritizes them. When it comes time to write up
findings in \texttt{RESEARCH.md}, the vocabulary from the brief ---
\texttt{MedicationStatement}, \texttt{Observation.code}, HIPAA,
\texttt{510(k)} --- shows up in the prose because the researcher saw
those terms as part of its working instructions.

\textbf{\large Step 4: Verify \texttt{RESEARCH.md}.}

Open \texttt{.planning/research/RESEARCH.md} and confirm the domain
knowledge landed. Here is an excerpt from the file our example
produced, covering the ``Drug data sources'' and ``Data format'' sections:

\begin{gsdplain}
## Drug data sources

Primary: FDA openFDA Drug Event API
  (https://api.fda.gov/drug/event.json). Use `search=` query
  parameter for adverse event lookups; rate limit 240 req/min
  unauthenticated.

Normalization: RxNorm via the NLM RxNav REST endpoint
  (https://rxnav.nlm.nih.gov/REST/). Use `rxcui` lookup to
  collapse brand names to concept IDs before interaction checks.

Interaction data: DrugBank academic tier. Cite DrugBank ID
  (format DBnnnnn) in all interaction records.

## FHIR Patient resource shape

- Patient.identifier  (array of Identifier)
- Patient.name        (array of HumanName)
- Patient.telecom     (array of ContactPoint, PHI)
- Patient.birthDate   (date)
- Patient.address     (array of Address, PHI)

PHI fields MUST be encrypted at rest per HIPAA Privacy Rule.
\end{gsdplain}

Twelve lines, and every one of them is grounded in something the
brief told the researcher to care about. Endpoints are cited with
real URLs. Resources are named by their exact FHIR field paths. The
HIPAA note is present because the brief flagged it. Compare this to
yesterday's output and the difference is not subtle --- this reads
like something a health-tech engineer would actually hand to a
coworker, not a generic tour of the internet.

\textbf{\large Step 5: Iterate.}

You will not get the brief perfect on the first try. Maybe the
researcher missed the SNOMED CT code system, or over-cited DrugBank
when RxNorm would have been more appropriate, or under-cited the
510(k) clearance pathway because your wording was too hedged. The
fix is straightforward: open \texttt{./skills/health-research/SKILL.md}
again, sharpen the instructions, save the file, and re-run
\texttt{/gsd:research-phase}. The next researcher spawn picks up the
edits without any rebuild step, because \texttt{agent\_skills} reads
live from disk on every spawn.

If you find that a single monolithic file is getting unwieldy, split
it into multiple markdown files inside the same directory. Name them
\texttt{01-sources.md}, \texttt{02-fhir.md}, \texttt{03-regulatory.md},
and GSD will concatenate them in lexicographic order. The subagent
still sees one continuous block, but you get to reason about the
pieces separately.

When you are happy with the brief, commit it. The
\texttt{./skills/health-research/} directory and the updated
\texttt{.planning/config.json} both belong in version control, and
every developer who clones the repo will get the same researcher
behavior automatically --- no setup steps, no onboarding friction,
no tribal knowledge about ``the way we do research on this project.''
The knowledge is written down, and it travels with the code.

\end{gsdscenario}

\begin{gsdtip}
Skill files are live. Edit and re-run --- there is no rebuild step.
If a GSD subagent is missing context, the fastest path to a fix is
almost always ``write it into the skill file and spawn again.''
\end{gsdtip}

\begin{gsdsuccess}
You have now seen the full shape of GSD's skills system: global skills
that live in your per-user directory and apply everywhere, legacy slash
commands for older Claude Code versions, twelve runtime targets that
all converge on the same user experience, and project-local
\texttt{agent\_skills} that let you teach GSD subagents anything your
domain needs them to know. Part II picks up from here and looks at
agents and teams --- the mechanism by which GSD composes subagents into
larger, coordinated work units.
\end{gsdsuccess}
% ============================================================================
%  part2-mcp.tex
%  Part II of "GSD: Going Further" --- MCP Integration
%  Chapters 5--9.
% ============================================================================

\part{MCP Integration}

\chapter{MCP Fundamentals for GSD Users}

If Part I was about teaching GSD to speak your project's dialect, Part II
is about teaching GSD to reach outside of itself. Most of the real work
developers do involves tools: design tools, databases, ticket systems,
cloud consoles, monitoring dashboards, that one custom script your team
lead wrote three years ago that nobody fully understands. Claude Code can
talk to all of them --- but only if a bridge exists. That bridge is MCP.

\section{What MCP Actually Is}

MCP --- the Model Context Protocol --- is an open standard, published by
Anthropic, for connecting AI assistants to external tools and data
sources. At its heart it is a small, stable JSON-RPC interface. A program
that implements MCP is called an \emph{MCP server}. A program that
connects to one is called an \emph{MCP client}. Claude Code is a client.
Almost everything else in this chapter is a server.

Each MCP server exposes a handful of \emph{tools}. A tool has a name, a
JSON schema describing its arguments, and a function that runs when Claude
Code invokes it. During a conversation, Claude Code sends a list-tools
request to every connected server at startup, merges the results into its
own tool surface, and from that point on, calling an MCP tool looks
identical to calling a built-in one. The assistant does not know or care
whether \texttt{read\_file} is implemented in Claude Code itself or served
by a filesystem MCP running on your laptop.

This is a much bigger deal than it sounds. Before MCP, every integration
between an assistant and an outside service had to be hand-coded, vendor
by vendor, with a bespoke protocol each time. MCP collapses that sprawl
into one JSON-RPC shape that anyone can implement in a weekend. You write
the server once, and every MCP-aware client --- Claude Code, Claude
Desktop, the various editor plugins --- can use it without modification.

\section{Why MCP Matters for GSD Specifically}

Here is the part that surprises people: you don't have to tell GSD
anything about your MCP servers. It figures them out on its own.

As of GSD v1.30 (see CHANGELOG entry \#1603), the executor and planner
subagents are explicitly instructed, in their system prompts, to inspect
the current set of available tools before drafting a plan or executing a
task, and to use any MCP tools they find as if they were first-class
capabilities. That was a one-line change to the subagent prompt templates
and it unlocked an enormous amount of behavior. If you connect a Postgres
MCP server to Claude Code, the planner now sees \texttt{postgres.query}
in its tool list and will naturally reach for it when you ask it to plan
a database migration. If you connect Google Stitch's MCP, the planner
sees \texttt{stitch.generate\_screen\_from\_text} and will write plan
steps that call it directly.

You don't edit a GSD config file to declare the tool. You don't hand-wire
anything in \texttt{.planning/}. The subagents inherit Claude Code's
entire tool surface, MCP and all, and they use what they find.

\begin{gsdtip}[You don't need to tell GSD about your MCP servers]
GSD's agents discover MCP tools automatically --- that's what v1.30's
\#1603 fix added. Just configure the MCP server in Claude Code and GSD
will pick it up on the next planning or execution run.
\end{gsdtip}

\section{Transport Types: stdio and HTTP}

MCP servers come in two flavors, distinguished by how the client talks
to them.

\subsection{stdio Transport}

The most common shape is \emph{stdio}: the client launches the server as
a local subprocess and exchanges JSON-RPC messages over the child's
standard input and standard output streams. This is simple, fast, and
completely local --- no network stack, no authentication, no ports to
open. Most of the MCP servers you will install for personal use are
stdio servers, invoked as \texttt{npx} or \texttt{uvx} commands. The
filesystem server, the GitHub server, the official Postgres reference
server, the Stitch proxy we'll meet in Chapter 7: all stdio.

The lifecycle is straightforward. When you start a Claude Code session,
the client reads your MCP config, spawns each stdio server as a
subprocess, speaks JSON-RPC to it over pipes, and keeps the subprocess
alive for the duration of the session. When the session ends, the
subprocess is killed. Restarting Claude Code restarts all your stdio
servers from a clean slate.

\subsection{HTTP Transport (Streamable HTTP)}

The other shape is remote HTTP, now formally called Streamable HTTP in
the MCP specification. Instead of spawning a subprocess, the client
connects to an HTTP endpoint over the network and exchanges messages as
HTTP requests with server-sent events for streaming. This is how you
would connect to a shared team MCP server, a cloud service that exposes
an MCP endpoint, or any tool that doesn't want to run on your machine.

HTTP transport is more complicated in exchange for more power: you have
to think about authentication (API keys in headers, bearer tokens, OAuth
flows), endpoint availability, and the fact that your network is now in
the critical path of every tool call. It is the right choice for
production integrations and the wrong choice for personal utilities.

\section{How Discovery Works}

When Claude Code starts a session, it walks a short, ordered list of
configuration sources and merges the results. The project-scoped file
\texttt{.mcp.json} in the current working directory comes first; the
user-scoped file \texttt{\~{}/.claude.json} comes second. A server name
defined in the project file takes precedence over the same name in the
user file, which means you can ship a project-local override for a
server that you otherwise use everywhere.

For each server entry, Claude Code either launches a subprocess (stdio)
or opens a connection (HTTP), sends the MCP handshake, and asks for the
server's tool listing. Those tools are added to the set of things the
assistant can invoke in the current conversation. GSD's subagents, when
they are spawned by commands like \texttt{/gsd:plan-phase} or
\texttt{/gsd:execute-phase}, inherit the same tool surface from the
parent Claude Code process. There is no second discovery step, no second
config file, no second place for things to go wrong.

The upshot is that an MCP server you add once is immediately available
everywhere GSD does useful work: planning, execution, verification,
review. The subagents don't need to be taught about the new tool. They
just find it sitting on the shelf.

\section{A Short History and Why It Matters}

MCP is young. The specification was first published in late 2024, the
reference implementations appeared shortly after, and the ecosystem of
third-party servers has been expanding roughly monthly since then. In
practical terms this means three things. First, the protocol surface
is still small enough to read in an afternoon --- if you're curious
about what an MCP message looks like on the wire, the spec is not yet
so sprawling that you need a guide to navigate it. Second, new
servers are appearing constantly, and the server you want probably
exists (or will exist by next month) even if it doesn't show up on
the first page of search results today. Third, the protocol is
versioned and backwards-compatible in a way that means the config
patterns in this guide are unlikely to break under you in the near
term --- when the spec evolves, it does so by adding new capabilities
to the negotiation step rather than breaking old message shapes.

The reason any of this matters for a GSD user is that the protocol's
youth shows up in rough edges. Error messages from MCP servers are
sometimes less polished than the equivalent error from a mature
vendor SDK; documentation for less-common servers is sometimes
scattered across a README, a changelog, and one Discord thread; and
the first time you wire up a new server, expect to spend ten minutes
figuring out which environment variable the server's author chose to
name the API key. None of this is fatal. All of it is the cost of
being a year early. The payoff is that the glue code you write today
will still work in three years, because the protocol itself is
designed to outlive any one vendor's implementation.

\section{What MCP Is Not}

A brief list of things MCP is not, because misconceptions in this
space are expensive.

MCP is \textbf{not} a replacement for HTTP APIs. If you have a REST
API that you love and that works, MCP does not make it faster or
better; it wraps it. The value of MCP is not performance. It is
standardization. An MCP server in front of your REST API means any
MCP-aware client can use it without hand-coding the integration.
That's worth something if you have more than one client. It's
worth less if you only ever use one.

MCP is \textbf{not} an agent framework. The protocol says nothing
about how the assistant decides which tool to call, how to chain
tool calls together, how to handle errors, or how to plan
multi-step work. That's the assistant's job. MCP just moves bytes
between a client and a server. The intelligence lives on either
side of the wire, not in the wire itself.

MCP is \textbf{not} a permissions system. Giving a tool to Claude
Code via an MCP server grants the assistant whatever access the
underlying credentials provide. If your Postgres MCP server
connects with a superuser account, Claude Code gets superuser
access to the database, and so does every subagent GSD spawns.
Scope your credentials at the database, filesystem, or API level
--- that's where permissions belong --- and treat MCP as a
transport layer that will faithfully carry whatever authority you
give it.

MCP is \textbf{not} magic. The tools you install are the tools you
get. If no one has written an MCP server for the thing you want to
integrate with, you have two options: find a close-enough existing
server (a generic HTTP fetch server, a shell-command runner, a
scripting server) or write one yourself. Writing one yourself is
surprisingly cheap --- the reference TypeScript and Python SDKs
make a minimum-viable server about fifty lines of code --- but it
is still work, and no amount of AI enthusiasm will conjure a
server that doesn't exist yet.

\chapter{Configuring MCP Servers}

Now that you know what MCP is and why GSD cares, let's add one. There are
three ways to register a server with Claude Code: the \texttt{claude mcp}
CLI, a project-scoped \texttt{.mcp.json} file, and the user-scoped
\texttt{\~{}/.claude.json} file. All three produce the same runtime
behavior; they differ only in scope and in how much typing you do.

\section{Adding Servers via the CLI}

The \texttt{claude mcp add} command is the fastest way to register a
server without hand-editing JSON. Two examples cover the two transports:

\begin{gsdbash}
# Register a local stdio server (the Stitch proxy, which we'll meet
# in Chapter 7). The "--" separator marks the end of claude's own
# flags and the start of the subprocess command line.
claude mcp add stitch --transport stdio -- npx @_davideast/stitch-mcp proxy

# Register a remote HTTP server at user scope (-s user means every
# project on this machine will see it).
claude mcp add my-api --transport http https://api.example.com/mcp -s user
\end{gsdbash}

The CLI writes the entry to the appropriate config file for you ---
project or user scope, depending on the \texttt{-s} flag --- and from that
point on, the server is indistinguishable from one you added by editing
JSON directly. Use the CLI for one-shot experimentation; use JSON files
when you want to commit the config alongside the rest of your project
or when you want finer control over env-var injection.

\section{Project-Scoped Configuration: .mcp.json}

A \texttt{.mcp.json} file in the root of your project declares servers
that are specific to this project and nowhere else. The shape is dead
simple:

\begin{gsdjson}
{
  "mcpServers": {
    "filesystem": {
      "command": "npx",
      "args": ["-y", "@modelcontextprotocol/server-filesystem", "/Users/me/projects"]
    }
  }
}
\end{gsdjson}

The top-level \texttt{mcpServers} object holds one entry per server,
keyed by a short name of your choosing (this name becomes the prefix on
tool calls --- so \texttt{filesystem.read\_file} would be the fully
qualified tool name in the above example). The \texttt{command} is the
binary to launch, and \texttt{args} is the argv passed to it. That's the
minimum viable config.

\section{User-Scoped Configuration: \textasciitilde/.claude.json}

The user-scoped file lives at \texttt{\~{}/.claude.json} and has the
same shape, just in a different place. Entries here apply to every
project you open in Claude Code, unless a project-local entry with the
same server name overrides it. Use the user-scope file for personal
utilities: your filesystem server pointing at \texttt{\~{}/projects},
your GitHub server with your personal access token, whatever daily-driver
search MCP you've settled on.

\section{Project vs User Scope: How to Decide}

The rule of thumb is simple. If the tool is specific to this codebase ---
a database server configured with this app's connection string, a Stitch
project that holds this app's UI, a staging API endpoint you only ever
hit from this repo --- use \textbf{project scope}. If the tool is a
personal utility that you'd want in every session --- filesystem,
GitHub, your search engine of choice --- use \textbf{user scope}. When
in doubt, start at user scope and move it down to project scope only if
it turns out to need project-specific credentials or paths.

A slightly more nuanced version of the same rule: ask yourself what
would happen if a teammate cloned your repo and ran Claude Code in
it. Would they need this server to reproduce your work? If yes ---
the database MCP, the Stitch MCP pointing at a shared project, the
deployment API wrapper --- it belongs in project scope, committed
(minus credentials) so the next developer gets it for free. If no
--- your personal filesystem server, a search MCP pointing at your
own notes, a weather server you installed for fun --- it belongs in
user scope where it stays with you and doesn't clutter the shared
config. The tiebreaker, when both feel true, is whether the tool
holds state that is specific to this project. A Postgres MCP
holding this app's connection string clearly does. A GitHub MCP
holding your personal access token clearly does not, even if you
only ever use it in one repo at a time.

There's a third scope worth knowing about, even though most users
will never touch it directly: the session-local scope created by
\texttt{claude mcp add} without either \texttt{-s project} or
\texttt{-s user}. This adds a server for the duration of the current
Claude Code session only, in-memory, with no file changes. It is
useful for truly one-shot experiments --- you want to poke at a new
server for ten minutes without committing to it --- and it vanishes
when you exit the session. Treat it as a try-before-you-buy mode.

\section{Environment Variables in Config}

Hardcoding an API key into \texttt{.mcp.json} is a loaded gun. Even if
you never commit the file, it still sits on your disk in plaintext,
readable by any process running as your user, and any casual
\texttt{cat .mcp.json} will print it to a terminal that might be
scrolled back or screen-shared or logged by an operator tool. The
correct pattern is env-var interpolation: store the secret in your
shell environment and reference it by name in the config.

\begin{gsdjson}
{
  "mcpServers": {
    "github": {
      "command": "npx",
      "args": ["-y", "@modelcontextprotocol/server-github"],
      "env": {
        "GITHUB_PERSONAL_ACCESS_TOKEN": "${GITHUB_TOKEN}"
      }
    }
  }
}
\end{gsdjson}

\credentialwarning

The \texttt{env} block passes environment variables into the subprocess.
Values of the form \texttt{\$\{NAME\}} are expanded from your own
environment at the moment Claude Code spawns the server --- so the real
token lives in your shell profile (or a secrets manager, or a dotenv
file that is itself gitignored), not in the config file. If the variable
is unset, the substitution leaves an empty string and the server will
almost always fail with an auth error; that's the correct failure mode.

\begin{gsdimportant}[.mcp.json belongs in .gitignore]
Add \texttt{.mcp.json} to \texttt{.gitignore} whenever it carries API
keys, OAuth client secrets, or database connection strings --- even if
you're using env-var interpolation, because the shape of the file
itself can leak information about your infrastructure. The user-scope
file \texttt{\~{}/.claude.json} lives outside any repository so it is
safer for shared credentials by default, but you should still treat it
like an SSH key: tight file permissions, no copies in shared drives,
no pasting into chat tools. A credential that leaves your laptop is a
credential that has leaked, regardless of where it ended up.
\end{gsdimportant}

One more footgun worth flagging: some editors and IDEs offer to
``helpfully'' open \texttt{.mcp.json} in a JSON schema view that pings a
remote schema registry. Disable this for any config file that contains
internal endpoint URLs or server names you'd rather not publish to
someone else's telemetry. The GSD guide team has seen this happen once.
Once is enough.

\chapter{Scenario --- Google Stitch + GSD}

\begin{gsdscenario}[Design-to-Code: A 5-Screen App with Stitch and GSD]

Google Stitch is an AI-powered UI generation product built on Gemini 2.5
Pro. You describe a screen in plain English (or hand it a reference
image), and it produces production-ready frontend code in the framework
of your choice --- HTML/CSS, React, Vue, Angular, all with Tailwind
variants. For developers who can ship backends in their sleep but whose
last three side projects died because the UI looked like a 2003
wireframe, this is the exact tool you've been waiting for.

Stitch ships with an official MCP integration in the form of the
\texttt{@\_davideast/stitch-mcp} npm package. The package provides two
things: an interactive \texttt{init} wizard that handles the painful
Google Cloud setup, and an MCP proxy server that your Claude Code
client connects to. The proxy is the recommended entry point, and the
rest of this scenario walks through why.

The Stitch MCP server exposes the following tools:
\texttt{create\_project}, \texttt{generate\_screen\_from\_text},
\texttt{get\_screen}, \texttt{get\_screen\_code},
\texttt{get\_screen\_image}, \texttt{build\_site},
\texttt{list\_projects}, and \texttt{list\_screens}. You won't need to
invoke any of them by hand --- GSD's planner and executor subagents
will find them in the tool list and wire them into the plan for you
once you've completed setup.

\vspace{6pt}
\textbf{Step 1 --- Admit that the problem is real.}
\vspace{4pt}

You have a meal-planning app in your head. Onboarding screen, dashboard,
recipe browser, meal calendar, shopping list: five screens total. You
can write the backend this weekend, and you have written it in your
head three times already. What has always stopped you is that the UI
from each of those mental iterations looked like a homework assignment.
You don't want to learn Figma. You don't want to be tied to one design
tool for the rest of your career. You do want to ship. This is the
problem Stitch is trying to solve, and the question this chapter
answers is what happens when you hand that solved problem to GSD.

\vspace{6pt}
\textbf{Step 2 --- Set up a Google Cloud project and enable the Stitch API.}
\vspace{4pt}

Open the Google Cloud Console (console.cloud.google.com), sign in with
the Google account you want to use for Stitch billing, and either
create a new project or pick an existing one from the project
dropdown. From the API \& Services section, search the product catalog
for ``Stitch'' and click Enable on the product page. You'll be asked
to confirm or attach a billing account; API usage bills against that
account at whatever rate Google currently publishes, so check the
current pricing page before you enable it if you're cost-sensitive.

Two notes about this step that will save you a headache. First, the
Stitch API is not enabled by default on any project, including new
ones --- you always have to click Enable explicitly. Second, an active
billing account is required; the free tier is metered, not permanent,
and a project without billing will fail at the first tool call with an
authentication-like error that isn't actually an authentication
problem. Make sure billing is attached.

\vspace{6pt}
\textbf{Step 3 --- Install stitch-mcp.}
\vspace{4pt}

\begin{gsdbash}
# Run the interactive initializer. It will walk you through
# gcloud auth, project selection, and config generation.
npx @_davideast/stitch-mcp init
\end{gsdbash}

The \texttt{init} wizard is the part you will love. It runs
\texttt{gcloud auth application-default login} for you if gcloud is
installed and you aren't already signed in, shows you the list of
projects you have Stitch access to, lets you pick one, and then prints
a ready-to-paste \texttt{.mcp.json} snippet at the end. It also tells
you which environment variables it expects for the API-key path if
you'd rather go that route instead of OAuth.

\vspace{6pt}
\textbf{Step 4 --- Choose an auth method.}
\vspace{4pt}

There are three ways to authenticate a Stitch MCP connection, and
getting this decision right on day one is worth more than any other
detail in this walkthrough. Here's the comparison:

\begin{center}
\begin{tabularx}{\linewidth}{@{}lXlX@{}}
\toprule
\textbf{Method} & \textbf{Setup} & \textbf{Token Lifetime} & \textbf{Best For} \\
\midrule
API Key & Set \texttt{STITCH\_API\_KEY} env var & Permanent & Solo developers, quick local testing \\
\addlinespace
OAuth (gcloud) & \texttt{gcloud auth application-default login} & 1 hour, auto-refreshed by proxy & Teams, CI/CD, daily driver \\
\addlinespace
Remote MCP (Google native) & \texttt{claude mcp add stitch --transport http} & 1 hour, MANUAL refresh & Quick experiments only \\
\bottomrule
\end{tabularx}
\end{center}

The numbers in the middle column are load-bearing. Google OAuth
application-default tokens have a \textbf{one-hour lifetime} --- this
is a Google-wide policy, not a Stitch-specific limit, and no amount of
configuration will extend it. What differs between rows 2 and 3 is who
handles the refresh. The local proxy from \texttt{@\_davideast/stitch-mcp}
watches the token's expiration, re-runs the refresh flow in the
background when the token gets close to expiring, and transparently
keeps your session alive for as long as Claude Code is running. A
manually-added remote MCP does not do this --- the token expires in
the middle of a working session, your next \texttt{generate\_screen}
call fails with a 401, and you have to manually re-authenticate before
you can continue.

For anything resembling daily work, the OAuth-via-proxy path is the
correct choice. The auto-refresh is the single biggest quality-of-life
feature in the whole integration. The API-key path is legitimately
fine for solo developers who want to get started quickly, but every
committed key is a potential leak and every rotation is manual; budget
for that if you go that way. The remote-MCP path is genuinely only
useful for a quick experiment: spin it up, poke at one screen, tear it
down. Anything longer and you'll spend more time re-authenticating
than building.

\vspace{6pt}
\textbf{Step 5 --- Add MCP config to \texttt{.mcp.json}.}
\vspace{4pt}

Paste the config block from the \texttt{init} wizard into a
\texttt{.mcp.json} file at the root of your project. It will look like
this (the proxy path, which is what you want):

\begin{gsdjson}
{
  "mcpServers": {
    "stitch": {
      "command": "npx",
      "args": ["@_davideast/stitch-mcp", "proxy"]
    }
  }
}
\end{gsdjson}

\credentialwarning

If you decided during Step 4 to go the API-key route instead, the
config shape is slightly different --- you supply the key through the
\texttt{env} block rather than relying on ambient gcloud credentials:

\begin{gsdjson}
{
  "mcpServers": {
    "stitch": {
      "command": "npx",
      "args": ["@_davideast/stitch-mcp", "proxy"],
      "env": {
        "STITCH_API_KEY": "${STITCH_API_KEY}"
      }
    }
  }
}
\end{gsdjson}

\credentialwarning

Reminder from Chapter 6: any \texttt{.mcp.json} that carries
credentials --- API keys in \texttt{env}, connection strings, bearer
tokens --- belongs in \texttt{.gitignore}. The env-var interpolation
pattern (\texttt{\$\{STITCH\_API\_KEY\}}) pushes the actual secret out
of the committed file and into your shell environment, which is
strictly better than a literal, but doesn't replace gitignoring the
file entirely for credential-bearing configs.

\vspace{6pt}
\textbf{Step 6 --- \texttt{/gsd:new-project}.}
\vspace{4pt}

Now the fun part. Fire up Claude Code in your empty project directory
and run:

\begin{gsdbash}
/gsd:new-project a meal-planning app with onboarding, dashboard,
  recipe browser, meal calendar, and shopping list --- 5 screens
  total. Use Stitch for UI generation with React and Tailwind.
\end{gsdbash}

GSD's new-project command writes \texttt{.planning/PROJECT.md} with a
one-paragraph description, a rough feature list, and a detected tech
stack. This is also the moment it kicks off the \emph{discuss} phase,
which is the structured conversation where design decisions get
locked in before any code is written. At this point GSD does not yet
know it will be using Stitch --- it has inferred ``React'' and
``Tailwind'' from your prompt, but it hasn't yet seen the Stitch tools
in its tool list because the discuss phase hasn't run. That's fine.
The next step is where the MCP layer starts to earn its keep.

\vspace{6pt}
\textbf{Step 7 --- \texttt{/gsd:discuss-phase 1}.}
\vspace{4pt}

The discuss phase is where you lock in decisions that the planner
will treat as constraints. This is where you commit to React +
Tailwind as the framework (so the planner doesn't hedge later), where
you decide whether to let Stitch pick a color palette or supply one
yourself, and where you give the planner the Stitch project ID it
will need to call the MCP tools. Here's a fragment of what you
might see in the resulting \texttt{DISCUSS.md}:

\begin{gsdyaml}
## Design Decisions --- Phase 1

Framework: React 18 + Tailwind CSS 3
UI tool: Google Stitch (via stitch-mcp proxy)
Stitch project ID: YOUR_PROJECT_ID
Color palette: defer to Stitch (will regenerate if needed)
Target screens: onboarding, dashboard, recipes, calendar, shopping

Rationale: Stitch generates React + Tailwind natively, so no
transpilation step. Letting Stitch pick the palette on the first
pass gives us a starting aesthetic to react to --- easier than
designing from a blank page.
\end{gsdyaml}

The Stitch project ID is the key piece. Without it the planner can
reach the Stitch MCP's tool list but won't know which project to
operate in; with it, every subsequent tool call will target the
correct project automatically. Grab the ID from Stitch's project
listing (or from \texttt{stitch.list\_projects} if you'd rather ask
the MCP server itself) and paste it into the discuss phase output.

\vspace{6pt}
\textbf{Step 8 --- \texttt{/gsd:plan-phase 1}.}
\vspace{4pt}

With the discuss phase committed, run plan-phase. The planner
subagent spawns with access to Claude Code's full tool surface ---
which now includes every tool the Stitch MCP exposes. It reads
\texttt{DISCUSS.md}, sees that Stitch is the chosen UI tool, sees
\texttt{stitch.generate\_screen\_from\_text} sitting in its tool
list, and writes plan steps that call it directly.

Here's the flavor of what the resulting \texttt{PLAN.md} will look
like once the planner is done:

\begin{gsdplain}
Task 3: Generate onboarding screen via Stitch
  Tool:   stitch.generate_screen_from_text
  Inputs: {
    project_id: YOUR_PROJECT_ID,
    prompt:     "Onboarding screen for a meal planning app. Three
                 steps: welcome, dietary preferences, household size.
                 Warm, approachable tone. Tailwind utility classes.",
    framework:  "react-tailwind"
  }
  Output: HTML/CSS in src/screens/Onboarding.tsx
  Depends on: Task 1 (project scaffold), Task 2 (stitch auth check)
\end{gsdplain}

The planner generates one of these blocks per screen --- five screens,
five Stitch calls --- plus the usual scaffolding, routing, and
integration tasks. Notice that the planner is writing tool-specific
inputs with real Stitch-compatible field names (\texttt{project\_id},
\texttt{framework}, \texttt{prompt}). It knows to do that because the
MCP tool's schema tells it what arguments the tool accepts. Schemas
beat prose every time.

\vspace{6pt}
\textbf{Step 9 --- \texttt{/gsd:execute-phase 1} and
\texttt{/gsd:ui-review 1}.}
\vspace{4pt}

Execute-phase is where the rubber meets the road. The executor
subagent walks through the plan task by task, calling MCP tools for
the tasks that declare them, and writing real files on disk as the
tools return. You'll see progress output along the lines of:

\begin{gsdplain}
[Task 3/14] Generate onboarding screen via Stitch
  -> Calling stitch.generate_screen_from_text
  -> Received screen_id = scr_a8f20b1e
  -> Calling stitch.get_screen_code (framework=react-tailwind)
  -> Wrote src/screens/Onboarding.tsx (2.4 KB)
  -> Calling stitch.get_screen_image (for reference asset)
  -> Wrote docs/screens/onboarding.png (47 KB)
[Task 3/14] COMPLETE
\end{gsdplain}

Each screen produces a real file in \texttt{src/screens/} and a
reference PNG in \texttt{docs/screens/}. The reference PNGs are not
strictly necessary but they're worth keeping: they're what Stitch's
renderer produced for the same prompt, and later on they make it
much easier to reason about visual drift when you iterate on the
copy or the palette.

Once the executor has run through all five screens, run
\texttt{/gsd:ui-review 1}. The ui-review command audits the generated
UI against GSD's six-pillar UI standards: accessibility, usability,
consistency, information density, responsive behavior, and visual
polish. It will flag missing alt text, contrast failures, keyboard
traps, and the other quiet ways AI-generated UI can be technically
correct and practically broken. Stitch's output is generally good,
but it is not infallible, and ui-review is where you find out.

\credentialwarning

\vspace{4pt}
\end{gsdscenario}

\begin{gsdtip}[Framework choice belongs in the discuss phase]
Stitch can generate in at least six frameworks (HTML/CSS, React,
Vue, Angular, with or without Tailwind). Tell GSD which framework
you want during the \emph{discuss} phase, not the \emph{execute}
phase. The planner will then write framework-specific task inputs
that the executor can hand directly to \texttt{stitch.get\_screen\_code}
without a conversion step. Waiting until execute-time means the
planner has already committed to a structure it now has to rewrite.
\end{gsdtip}

\chapter{Scenario --- Database MCP + GSD Planning}

\begin{gsdscenario}[Schema-Aware Planning: Postgres MCP for Migration Plans]

GSD's planner is smart but not clairvoyant. Ask it to ``add a
\texttt{user\_preferences} table'' in a project where it has no
visibility into the database, and it will produce a plausible migration
against a plausible schema --- not your schema. The columns it picks
will be reasonable guesses, the foreign key names will look like what a
textbook would suggest, and the migration will fail the moment it meets
the real database. Worse, a junior reviewer glancing at the plan might
not notice.

A Postgres MCP server fixes this by giving the planner read-only
introspection into the actual schema. With the MCP installed, the
planner issues real \texttt{SELECT} statements against
\texttt{information\_schema} before it drafts the migration, and the
resulting plan references the columns that actually exist. The rest of
this chapter walks through the setup.

\vspace{6pt}
\textbf{Step 1 --- Install a Postgres MCP server.}
\vspace{4pt}

The Anthropic-maintained reference server
\texttt{@modelcontextprotocol/server-postgres} is the safest starting
point. It is deliberately \textbf{read-only}: the only tool it exposes
is a \texttt{query} function that runs against the database, and the
server wraps every query in a read-only transaction so even a
\texttt{DELETE} inside the query text cannot mutate data. That's a
feature, not a limitation. For planning work, read-only is all you
need and all you want.

\begin{gsdbash}
# No install step is needed -- npx fetches the package on demand
# when Claude Code launches the subprocess. But you can pre-cache
# it to shave a few seconds off the first session:
npx -y @modelcontextprotocol/server-postgres --version
\end{gsdbash}

\vspace{6pt}
\textbf{Step 2 --- Configure with a connection string.}
\vspace{4pt}

Add the server to your \texttt{.mcp.json} with a connection string
supplied as an environment variable. Never hardcode the string into
the file --- database URLs contain credentials by definition, and a
credential in a config file is a credential one \texttt{git add} away
from leaving the building.

\begin{gsdjson}
{
  "mcpServers": {
    "postgres": {
      "command": "npx",
      "args": ["-y", "@modelcontextprotocol/server-postgres"],
      "env": {
        "DATABASE_URL": "${DATABASE_URL}"
      }
    }
  }
}
\end{gsdjson}

\credentialwarning

Export \texttt{DATABASE\_URL} in your shell before starting Claude
Code (from a dotenv file, your secrets manager, or a line in your
shell profile --- whichever you prefer). The substitution happens
when Claude Code spawns the MCP subprocess, so the value needs to be
set in the shell that launches Claude Code, not just in Claude Code's
own session. And yes: this \texttt{.mcp.json} belongs in
\texttt{.gitignore} for exactly the reasons Chapter 6 covered.

\vspace{6pt}
\textbf{Step 3 --- \texttt{/gsd:plan-phase} with schema visibility.}
\vspace{4pt}

With the Postgres MCP running, start a plan-phase command that will
require schema knowledge. For example:

\begin{gsdbash}
/gsd:plan-phase 2 add a user_preferences table with notification
  toggles, theme selection, and default meal calendar view
\end{gsdbash}

The planner subagent sees \texttt{postgres.query} in its tool list
and does the right thing without being told: before drafting any
migration steps, it inspects the existing schema with queries like:

\begin{gsdplain}
SELECT column_name, data_type
FROM information_schema.columns
WHERE table_name = 'users'
ORDER BY ordinal_position;
\end{gsdplain}

It does this for every table the new migration might touch, not just
the obvious ones, because it was trained to build context before
acting. The returned rows get folded into the planning context, and
the migration steps the planner writes are grounded in what actually
exists.

\vspace{6pt}
\textbf{Step 4 --- Observe \texttt{PLAN.md} referencing real columns.}
\vspace{4pt}

The resulting plan will include a block that looks like this:

\begin{gsdplain}
Task 5: Migration -- create user_preferences table
  Preconditions verified via postgres.query:
    - users table exists with columns (id, email, created_at,
      updated_at, last_login_at)
    - no existing user_preferences table
  Migration SQL:
    ALTER TABLE users ADD COLUMN preferences JSONB;
    -- or, if we prefer a separate table:
    CREATE TABLE user_preferences (
      user_id INTEGER PRIMARY KEY REFERENCES users(id),
      ...
    );
\end{gsdplain}

That comment block listing preconditions is the key output of the
whole exercise. A human reviewer can look at the plan and verify, at
a glance, that the planner saw the real schema when drafting the
migration. The reviewer's job is no longer ``spot-check every column
name against the real database'' --- it is the much smaller job of
``confirm the preconditions the planner listed are the right ones.''

\vspace{6pt}
\textbf{Step 5 --- Cleanup.}
\vspace{4pt}

There is no cleanup. The Postgres MCP server is read-only and
stateless: it holds no data on disk, writes no files, and makes no
changes to your database. Killing Claude Code ends the subprocess;
starting a new session respawns it from scratch with the same config.
The only persistent trace is whatever lives in your query logs, which
is a database-side concern and not a GSD one.

\vspace{4pt}
\end{gsdscenario}

\begin{gsdimportant}[Read-only means read-only]
The Postgres MCP server is read-only. If you want the generated
migrations to actually run, that's a separate concern --- GSD plans
them, you (or your CI) execute them through your normal migration
tool (Flyway, Liquibase, Alembic, sqlx-migrate, whatever you already
use). Do not look for an MCP server that will run migrations for
you. The separation of planning from execution here is not a
limitation; it is the safety boundary that keeps a planning error
from becoming a production outage.
\end{gsdimportant}

\chapter{Building MCP-Aware GSD Workflows}

The two scenarios in Chapters 7 and 8 each used one MCP server in
isolation. Real projects accumulate servers the way shell profiles
accumulate aliases --- a filesystem server, a GitHub server, a
Postgres server for the main database, maybe a separate analytics
Postgres, Stitch for UI, a custom server that wraps your internal
deployment API. This chapter is about making the whole pile work
together, and about the edges where MCP integration gets
interesting.

\section{Combining Agent Skills with MCP}

Part I of this guide introduced agent skills --- the small project-
scoped prompt fragments that get injected into a subagent's context
whenever it runs in this project. Skills and MCP tools compose in a
way that is more powerful than either one alone.

The trick is that a skill can refer to an MCP tool by name. A skill
is just prose injected into the planner's or executor's prompt, and
if that prose says ``when planning database migrations, always call
\texttt{postgres.query} first to inspect existing tables,'' the
planner will do exactly that. The skill doesn't need to know how
\texttt{postgres.query} works; it just needs to mention the name.
The MCP server already told Claude Code about the tool's schema
during discovery, so the planner can fill in the arguments correctly
on its own.

This pattern is how you bake team conventions into a project. Your
team decides that every new feature needs a database schema check
before planning? Write a one-paragraph skill that says so, mention
the \texttt{postgres.query} tool by name, and every subsequent
\texttt{/gsd:plan-phase} run will follow the rule without anyone
having to remember to invoke it. The same trick works for any MCP
server: a skill that says ``for any UI task, always call
\texttt{stitch.list\_projects} first to confirm the Stitch project
exists'' costs you a paragraph and saves you a class of
``undefined project'' errors you would otherwise hit at execute
time.

A worked example makes this concrete. Suppose your team has agreed
that every database-touching phase must start by reading the current
schema, and every UI-touching phase must start by confirming the
Stitch project exists and has the expected screen count. You can
bake both rules into a skill like this (placed in
\texttt{.claude/skills/team-conventions.md} or wherever your
project's agent skills live):

\begin{gsdplain}
# Team conventions for planning and execution

## Before any database migration task
1. Call postgres.query with a schema inspection query
   against information_schema.tables and information_schema.columns
   for the affected tables.
2. Include the returned column list in the task preconditions.
3. Only then draft the migration SQL.

## Before any Stitch-generated UI task
1. Call stitch.list_projects and confirm the configured project
   ID is present.
2. Call stitch.list_screens for that project and record the
   current screen count.
3. If the screen count differs from what DISCUSS.md claims, stop
   and surface the mismatch instead of generating blindly.
\end{gsdplain}

The skill is maybe thirty lines of prose. The behavior change is
substantial: you stop finding out about schema drift at execute
time and start finding out about it at plan time, when fixing it
is cheap. Notice that this pattern does not require any changes to
GSD's source code, to the MCP servers themselves, or to Claude
Code's configuration. It is entirely text, and the text flows
through the same discovery-and-injection paths you already use for
every other skill.

The reason this works is worth pausing on. Subagent prompts are
assembled fresh at the start of each subagent invocation: the base
template, the current phase's context, any applicable skills, and
the current tool manifest. When a skill mentions a tool name, the
planner sees that mention \emph{and} the tool's schema side by
side. It doesn't have to guess what arguments the tool takes,
because the schema is right there. It doesn't have to be reminded
to use the tool, because the skill says so in plain language.
Skills and MCP compose cleanly because they're both contributing
to the same assembled prompt from different directions.

\section{The GSD Plugin Format (Issue \#1883)}

There's a larger integration idea in the pipeline. GSD issue \#1883
tracks a plan to package GSD itself as a Claude Code plugin --- a
single installable unit that would bundle GSD's skills, subagent
templates, and a dedicated MCP server, all in one place. The
promise is that a user could run one install command and get the
entire GSD experience, with its own MCP surface exposed as a first-
class tool set, rather than having GSD live alongside MCP as a
separate concern.

This is on the roadmap for v1.32 or later --- the plugin format is
still evolving, and the schema the plugin manifest will use is not
yet fixed. Committing to a specific shape in this guide would be
premature. The thing to know is that it's coming, that it is
specifically designed to make MCP integration a first-class part of
GSD rather than an afterthought, and that the scenarios in this
part will still work on the plugin version because the underlying
primitives are unchanged. Follow the issue if you want early visibility.

\section{Failure Modes: When an MCP Server is Unavailable}

An MCP server can fail mid-execution. The subprocess crashes, the
remote endpoint returns a 500, the network is flaky, the auth token
expires at exactly the wrong moment. What happens then?

Claude Code returns a tool error to the agent that made the call.
The error is not silent: it propagates back as a failed tool
response with an error message attached. GSD's executor subagent
sees the error and does one of two things.

The first option is a retry, chosen when the failure looks
transient --- a network timeout, a 503, an obviously flaky response.
The executor will retry the same tool call, sometimes with a short
delay, before concluding that the problem is real. This handles
most of the everyday flakiness you encounter with remote servers.

The second option is to mark the task as blocked and write a
deviation note to the execution log. A deviation note is the
executor's way of saying ``I can't finish this task and the reason
is not my fault.'' It includes the original plan step, the tool
error, and enough context for a human to resume work later. The
phase does not silently complete; it ends in a blocked state, and
the status command (\texttt{/gsd:status} or similar) will surface
the blockage the next time you ask.

The plan does not degrade silently. Failed MCP calls are visible in
the executor log, deviation notes make them unmissable, and the
phase-completion gate refuses to pass a phase where execution was
blocked on a tool error. If you've been running GSD long enough to
trust it, this is the behavior you want --- the tool doesn't lie
about what happened, it tells you something broke and stops.

\section{Observability: Knowing What Your MCP Servers Did}

An underappreciated part of running MCP at scale is knowing what
your servers have actually been doing. Claude Code's conversation
log shows you every tool call the assistant made, but once you have
five or six MCP servers running, that log can get noisy, and
correlating ``this file on disk'' with ``this tool call that
produced it'' is not always easy from the transcript alone.

Three patterns help. First, most MCP servers can be started in a
verbose or debug mode that prints each tool invocation to stderr.
Claude Code surfaces server stderr in its own logs, so turning on
verbose mode on a suspect server is often the fastest way to catch
a misbehaving call. The flag is server-specific (look for
\texttt{--verbose}, \texttt{--log-level debug}, or an environment
variable like \texttt{LOG\_LEVEL=debug}), but it exists almost
everywhere because MCP server authors are themselves debugging
their own servers against Claude Code and need the same signal.

Second, for servers you write or control, emit a one-line audit
record per tool call: timestamp, tool name, argument summary, outcome,
duration. Pipe it to a file or to your logging pipeline. When
something goes sideways at execute time, the audit log tells you
what actually happened versus what the plan says was supposed to
happen --- and those two narratives are not always the same.

Third, for production-grade setups, consider wrapping your MCP
servers behind a single reverse-proxy MCP that logs all traffic
and then forwards calls to the appropriate underlying server. This
adds a hop but gives you one place to look for everything. It's
overkill for personal projects and exactly right for shared team
infrastructure where ``who called what, when, against which
database'' is a compliance question with an actual answer.

\section{Version Drift: What Happens When a Server Changes}

Your MCP servers will change under you. The Stitch proxy will
publish a new version that renames a tool argument. The Postgres
server will add a new capability you've been wanting. A third-party
server you depend on will be deprecated in favor of a successor
that is almost but not quite backwards-compatible. None of this is
exotic; it's normal package ecosystem behavior, and MCP is no
different from npm or pip in that regard.

GSD copes with this better than you might expect, because the
discovery step happens at the start of every session. The planner
sees the current set of tools with the current schemas --- not a
cached snapshot from when you first configured the server. If a
tool's argument list changes overnight, the next plan-phase run
will use the new shape without you having to update anything in
GSD. If a tool disappears entirely, the planner will stop proposing
it, because it won't see it.

The failure mode to watch for is a skill that names a tool
explicitly. If your team-conventions skill says ``always call
\texttt{postgres.query}'' and someone upgrades the Postgres MCP to
a version where that tool was renamed to \texttt{postgres.execute},
the skill will keep asking for the old name, the planner will try
to call it, and the call will fail. The fix is to update the skill
to match. A short CI check that grep-s your skills for known MCP
tool names and cross-references them against a list of currently-
installed servers is cheap and will save you a debug session.

\begin{gsdnote}[Best practices summary]
A short checklist of what you should take away from Part II:
\begin{itemize}
  \item \textbf{Project-scoped servers} for project-specific tools: the
        database for this app, the Stitch project for this UI, the
        staging endpoints for this service.
  \item \textbf{User-scoped servers} for personal utilities: filesystem,
        GitHub, your daily-driver search MCP, anything you'd want
        available in every Claude Code session.
  \item \textbf{Env-var your credentials} every time, without exception.
        Never paste a real API key into a committed file, and never
        paste one into a file you're ``about to'' gitignore.
  \item \textbf{\texttt{.mcp.json} belongs in \texttt{.gitignore}} any
        time it carries secrets, even if you're using env-var
        interpolation. The shape of the file itself can leak
        information you don't want to publish.
  \item \textbf{For teams, prefer OAuth with service accounts} over
        shared API keys. Per-user auth with proper rotation is the
        only approach that survives a team member leaving. Shared
        keys are a time bomb with a known fuse length.
  \item \textbf{Token expiration is real}: OAuth tokens die after an
        hour. Use proxies that auto-refresh. Do not use raw remote
        MCP auth for anything you plan to run longer than a
        demo.
\end{itemize}
\end{gsdnote}

The overarching theme of Part II is that MCP is plumbing, and good
plumbing disappears. Once your servers are configured, the
gitignore is in place, and the env vars are exported, you stop
thinking about MCP at all. GSD's planner writes plans that use the
tools you've plugged in, the executor calls them, and the outputs
land on disk. The configuration work is upfront and small. The
payoff is that every subsequent plan and every subsequent execution
has access to the real state of your project --- the real schema,
the real Stitch project, the real repository --- and the quality of
the resulting work rises accordingly. That is the entire argument
for connecting GSD to MCP, and it is the reason Part III will
assume all of this is in place when it turns to multi-repo and
monorepo workflows.
% ============================================================================
%  part3-multirepo.tex
%  Part III of "GSD: Going Further" --- Multi-Repo Workflows
%  Chapters 10--14.
% ============================================================================

\part{Multi-Repo Workflows}

\chapter{Workstreams vs. Workspaces --- When to Use Which}

Parts I and II taught GSD to speak your project's dialect and to reach
outside of itself. Part III is about a more mundane but equally painful
problem: real projects are rarely a single repository with a single
timeline. Even when a project \emph{is} a single repo, it often runs
several independent tracks of work in parallel --- a backend feature,
a frontend redesign, an infrastructure migration, a one-off experiment
--- each with its own planning cadence, its own owners, and its own
definition of ``done''. And when the project genuinely spans multiple
repositories, as microservice fleets and polyrepo libraries often do,
the coordination problem multiplies.

GSD has two different answers to that problem, and learning when to
reach for each is the difference between a smooth multi-track workflow
and a pile of overwritten STATE.md files.

\section{The Problem: Shared Planning State}

The default assumption baked into GSD is that a repository has one
active plan at a time. The files in \texttt{.planning/} --- STATE.md,
ROADMAP.md, REQUIREMENTS.md, the phase directories --- describe a
single timeline of intent. Commands like \texttt{/gsd:progress} and
\texttt{/gsd:next} read that timeline to decide what to tell you.
\texttt{/gsd:execute-phase} writes to it.

If you try to run two unrelated workstreams against that shared state
at the same time, bad things happen. STATE.md flips between them as
each track edits it. The phase numbering collides. \texttt{/gsd:progress}
reports ``phase 4 in progress'' when you meant that for the backend
feature, and the frontend lead reads it and wonders why their redesign
lost two phases overnight. The ROADMAP gets clobbered because the
last writer wins. Every single one of those symptoms traces back to
the same root cause: two independent timelines trying to share one
set of planning files.

Workstreams and workspaces both solve that root cause. They solve it
at different layers of the stack, and they have different costs.

\section{Workstreams: One Repo, Many Timelines}

Workstreams were introduced in GSD v1.28. The idea is dead simple:
\emph{if you already have one repo, keep one repo, but give each
independent timeline its own isolated planning subtree.} Under the
hood, a workstream is a directory under \texttt{.planning/workstreams/}
that contains its own STATE.md, ROADMAP.md, REQUIREMENTS.md, and
phase directories --- a complete parallel planning tree, scoped
under a name of your choice. When you switch to a workstream, GSD
transparently re-roots all its planning operations to read and write
from that subtree. The root \texttt{.planning/} is untouched.

The commands are short and predictable. \texttt{/gsd:workstreams create
<name>} creates a new one. \texttt{/gsd:workstreams switch <name>}
makes it the active workstream for subsequent commands. \texttt{/gsd:workstreams
list} shows all of them. \texttt{/gsd:workstreams complete <name>} archives
one when its feature ships. The switch itself is essentially free ---
GSD just updates a pointer to the active workstream directory --- so
you can flip between two or three of them many times a day without
thinking about the cost.

What workstreams \emph{don't} do is give you git isolation. The code
is still the same checkout. If two workstreams both want to edit
\texttt{src/billing.ts} on the same branch, they will conflict in the
same way they would without workstreams. Workstreams isolate
\emph{planning state}, not \emph{source state}. For many real projects
that distinction is exactly right, because the planning collision is
the painful one and the source collision is already handled by git.

\section{Workspaces: Many Repos, One Plan}

Workspaces are the heavier instrument, also introduced in v1.28. A
workspace is a separate directory --- by convention under
\texttt{\~{}/gsd-workspaces/<name>/} --- that contains one or more
named repositories, each checked out as either a git worktree or a
full clone, plus a workspace-level manifest file called WORKSPACE.md.
Each repo inside the workspace has its own independent
\texttt{.planning/} subtree. The workspace itself is the coordination
layer: it knows which repos are in the workspace, tracks cross-repo
status in its manifest, and lets you run coordinated planning without
having to merge the repos into a monorepo you didn't want.

The commands mirror the workstream ones but operate at the repo-set
level. \texttt{/gsd:new-workspace --name v2-launch --repos api,web-frontend,admin-dashboard}
creates a workspace named \texttt{v2-launch} containing three repos.
\texttt{/gsd:list-workspaces} shows all configured workspaces.
\texttt{/gsd:remove-workspace v2-launch} tears one down when it's done.
The \texttt{--strategy} flag takes one of two values:

\begin{itemize}
\item \texttt{worktree} (the default) creates a git worktree of each
  named repo. Worktrees share object storage with the parent
  checkout, so they're cheap to create and cheap to keep around.
  The trade-off is that you can't have two worktrees on the same
  branch of the same repo at the same time --- git refuses.
\item \texttt{clone} performs a full \texttt{git clone} of each
  named repo into the workspace directory. This is slower and uses
  more disk, but it gives you complete isolation --- the workspace's
  copy of each repo is genuinely independent from your main checkout
  and can sit on any branch without conflict.
\end{itemize}

Workspaces cost more to create than workstreams because they involve
actual git operations --- worktree creation or cloning --- and a new
directory tree on disk. A workspace with three repos takes a few
seconds to materialize, not the sub-millisecond switch of a workstream.
That's the right cost when you actually need repo-level isolation,
and the wrong cost when you don't.

\section{Decision Table}

The shortest path to the right answer is to ask: \emph{do my parallel
timelines share code, or do they share planning?} If they share code,
a workstream is what you want. If they need different code or different
git state, a workspace is what you want.

\begin{center}
\small
\begin{tabularx}{\linewidth}{@{}lXX@{}}
\toprule
\textbf{Scenario} & \textbf{Use} & \textbf{Why} \\
\midrule
Backend + frontend in same repo &
Workstream &
Same code, different planning timelines. The extra cost of a
workspace buys you nothing. \\
\addlinespace
Three microservices needing coordinated planning &
Workspace (multi-repo) &
Different repos, each already has its own
\texttt{.planning/}. The workspace is the coordination layer. \\
\addlinespace
Feature branch isolation, current repo only &
Workspace \texttt{--repos .} &
A single-repo workspace gives you a worktree of the current repo
with its own independent planning state. \\
\addlinespace
Disposable spike or prototype &
Workspace \texttt{--strategy clone} &
Full isolation. When the spike dies, remove the workspace and
nothing leaks back into the main checkout. \\
\addlinespace
Refactor touching several repos atomically &
Workspace \texttt{--strategy worktree} &
Coordinated planning across the set, shared git fast-path under
the hood. \\
\bottomrule
\end{tabularx}
\end{center}

\begin{gsdtip}[Default to the cheaper answer]
Workstreams are much cheaper than workspaces. If the work you're
about to start fits in a single repo, default to creating a workstream
and only escalate to a workspace if you genuinely need git-level
isolation. It is always easier to graduate a workstream into a
workspace later than it is to undo a premature workspace.
\end{gsdtip}

\section{One More Consideration: How Long Does the Work Live?}

A secondary question worth asking before you pick is how long the
timeline is expected to live. Workstreams are well suited to long-lived,
repeating tracks: ``the backend team's current feature'', ``the ongoing
observability migration'', ``this quarter's performance work''. They
rotate through many features over their lifetime, and the fact that
their STATE.md persists across those features is useful.

Workspaces are better suited to \emph{cohorts of work that share a
release} --- the kind of thing where you'd otherwise say ``we're
shipping v2 across three services and I need to track all three
together''. When the cohort ships, you remove the workspace. The
repos themselves live on, but the workspace's WORKSPACE.md and its
coordination state get archived or deleted along with the launch
it tracked.

Neither pattern is wrong for the other's use case, but using a
workstream for a launch cohort tends to feel fiddly (you end up
making a workstream in each repo and hand-coordinating) and using
a workspace for a long-lived single-repo track tends to feel heavy
(you end up with a workspace directory that never quite gets cleaned
up). Match the lifetime of the construct to the lifetime of the work.

\chapter{Scenario --- Monorepo with Independent Apps}

\begin{gsdscenario}[NX-Style Monorepo: Backend, Frontend, and Mobile]

You run a single git repository containing three apps under
\texttt{apps/}: a Node backend API, a Next.js frontend, and a React
Native mobile app. The three teams work on different cadences ---
the backend team is mid-migration to a new ORM, the frontend team
is redesigning the onboarding flow, the mobile team is fighting a
long-tail of platform-specific bugs. They all share the same git
history and a fair amount of code under \texttt{packages/}, so
splitting into three repositories would be more disruptive than it
is worth. But running all three teams' GSD plans against the same
root \texttt{.planning/} has already produced one STATE.md corruption
this month, and you are done with that.

This scenario walks through giving each team its own workstream,
running an independent GSD flow in each, and archiving them cleanly
when features ship. Total elapsed time: about five minutes of
commands, most of which are the normal GSD flow you already know.

\textbf{Step 1. Starting layout.}

The repository looks like this before you touch anything:

\begin{gsdplain}
my-monorepo/
|-- apps/
|   |-- backend/
|   |-- frontend/
|   `-- mobile/
`-- .planning/
    |-- workstreams/
    |   |-- backend-api/
    |   |   |-- ROADMAP.md
    |   |   `-- STATE.md
    |   `-- frontend-redesign/
    |       |-- ROADMAP.md
    |       `-- STATE.md
    `-- config.json       <- shared config across workstreams
\end{gsdplain}

\noindent
(The \texttt{.planning/workstreams/} subtree is what we're about to
create. The root \texttt{.planning/config.json} already exists and
stays where it is --- it holds the shared config all workstreams
inherit from.)

\textbf{Step 2. Create the backend workstream.}

\begin{gsdbash}
/gsd:workstreams create backend-api
\end{gsdbash}

\noindent
GSD creates \texttt{.planning/workstreams/backend-api/} and
scaffolds a fresh set of planning files inside it --- an empty
ROADMAP.md, an initial STATE.md pointing at an un-planned phase,
and a small WORKSTREAM.md marker file that the tooling uses to
detect that the subtree is a workstream rather than a plain
directory. The root \texttt{.planning/} is not touched. You should
see output along the lines of:

\begin{gsdplain}
Created workstream: backend-api
Path: .planning/workstreams/backend-api/
Status: active
\end{gsdplain}

\textbf{Step 3. Switch to the new workstream.}

\begin{gsdbash}
/gsd:workstreams switch backend-api
\end{gsdbash}

\noindent
Switching makes \texttt{backend-api} the active workstream. Every
subsequent \texttt{/gsd:*} command in this session will read and
write state from \texttt{.planning/workstreams/backend-api/} instead
of the root \texttt{.planning/}. There's no second-command mystery
here: if you forget to switch, \texttt{/gsd:progress} will tell you
which workstream it's reporting on in its first line, so you catch
the mistake early.

\textbf{Step 4. Run a normal GSD flow, scoped to backend.}

At this point, everything looks and feels like vanilla GSD. You
start a new project for the ORM migration:

\begin{gsdbash}
/gsd:new-project
/gsd:discuss-phase 1
/gsd:plan-phase 1
/gsd:execute-phase 1
\end{gsdbash}

\noindent
The PLAN.md, STATE.md, and phase directories that come out of this
flow all live under \texttt{.planning/workstreams/backend-api/}.
The root \texttt{.planning/} does not see any of it. If you run
\texttt{/gsd:progress} right now, you'll get the backend workstream's
view --- phase 1, in progress, whatever your executor is chewing on.

\textbf{Step 5. Create the frontend workstream in parallel.}

While the backend team keeps going, the frontend lead wants to start
the redesign plan. In a new terminal (or just later in the same
terminal), they run:

\begin{gsdbash}
/gsd:workstreams create frontend-redesign
/gsd:workstreams switch frontend-redesign
/gsd:new-project
/gsd:discuss-phase 1
\end{gsdbash}

\noindent
This creates a second, entirely independent planning subtree at
\texttt{.planning/workstreams/frontend-redesign/}. The two workstreams
don't see each other's state files. They can be in completely
different phases, have completely different agent configurations,
use different \texttt{project\_code} prefixes --- the isolation is
at the planning-state level, not at the git level.

\textbf{Step 6. List workstreams.}

When you want to see what's going on across the whole repo, ask:

\begin{gsdbash}
/gsd:workstreams list
\end{gsdbash}

\noindent
The output is a small table, one row per workstream, giving you
name, status, current phase, and when it was last touched:

\begin{gsdplain}
NAME                STATUS    PHASE   LAST TOUCHED
backend-api         active    1       3 minutes ago
frontend-redesign   active    1       1 minute ago
\end{gsdplain}

\noindent
This is the one command that crosses workstream boundaries. Everything
else respects the active-workstream scope and only shows you one
workstream's view at a time.

\textbf{Step 7. Complete and archive.}

Two weeks later, the backend ORM migration ships. You finish
the last phase, run your tests, push the merge, and then retire
the workstream:

\begin{gsdbash}
/gsd:workstreams complete backend-api
\end{gsdbash}

\noindent
Completion does not delete anything. GSD moves the workstream's
directory to \texttt{.planning/workstreams/archive/backend-api/},
sets its status to \texttt{archived}, and stops surfacing it in
the default \texttt{list} output. If you ever need to resurrect it
--- say to investigate a post-ship question --- you can switch back
to the archived workstream and treat it as read-only evidence.

\textbf{Step 8. Cleanup.}

Repeat the \texttt{complete} step for the frontend workstream when
its redesign ships. The root \texttt{.planning/} has stayed clean
throughout: all the actual work lived under
\texttt{.planning/workstreams/}, and all the archived evidence lives
under \texttt{.planning/workstreams/archive/}. Your repo is ready
for the next set of parallel tracks, and there's a searchable record
of what the old tracks did.

\end{gsdscenario}

\begin{gsdtip}[Workstream scope is per-workstream]
Each workstream's \texttt{/gsd:progress} only reports on that
workstream's state. If you want a cross-workstream view, run
\texttt{/gsd:workstreams list}. This is the right scoping default
--- a backend dev should see backend state by default, not be
distracted by frontend state --- but it does mean that
``all-workstreams'' views are an explicit command rather than the
default.
\end{gsdtip}

\chapter{Scenario --- Microservice Fleet}

\begin{gsdscenario}[Three Repos, One Workspace: API + Web + Admin]

You have three separate repositories on disk: \texttt{api},
\texttt{web-frontend}, and \texttt{admin-dashboard}. They're
independent services with independent deployment pipelines, but
they all ship together as part of the v2 launch, and the product
manager would like a single view of ``how is v2 going?'' that
doesn't require opening three different planning files in three
different checkouts. A workspace is exactly the right shape for
this: it gives you a coordination layer on top of the three repos
without forcing you to merge them into a monorepo.

This scenario walks through creating the workspace, running
coordinated planning across all three repos, and tearing the
workspace down cleanly when v2 ships.

\textbf{Step 1. Create the workspace.}

\begin{gsdbash}
/gsd:new-workspace --name v2-launch \
  --repos api,web-frontend,admin-dashboard
\end{gsdbash}

\noindent
GSD looks up each named repo --- either via relative paths, or
via registered repo aliases (you can list the registered aliases
with \texttt{/gsd:list-workspaces}'s config output) --- and
materializes the workspace directory under
\texttt{\~{}/gsd-workspaces/v2-launch/}. For each repo, GSD creates
a git worktree (the default strategy) pointing at the repo's
current branch and drops it inside the workspace. It also writes
a WORKSPACE.md manifest at the workspace root listing the repos
and their current phase status.

The whole operation takes a few seconds. At the end of it, you
have a fully usable workspace containing three independently
planned repositories, sitting side by side.

\textbf{Step 2. Inspect the workspace tree.}

\begin{gsdplain}
~/gsd-workspaces/v2-launch/
|-- WORKSPACE.md          <- workspace manifest (orchestrator-owned)
|-- api/                  <- repo 1 (independent .planning/)
|   |-- .planning/
|   |   |-- ROADMAP.md
|   |   `-- STATE.md
|   `-- src/
|-- web-frontend/         <- repo 2 (independent .planning/)
|   |-- .planning/
|   `-- src/
`-- admin-dashboard/      <- repo 3 (independent .planning/)
    |-- .planning/
    `-- src/
\end{gsdplain}

\noindent
WORKSPACE.md is the orchestrator-owned manifest --- it tracks which
repos belong to the workspace and what the cross-repo phase status
is. Each repo below it is an independent checkout with its own
\texttt{.planning/} subtree. The three \texttt{.planning/} trees
are genuinely independent: they don't share state, they don't share
phase numbers, they don't share REQUIREMENTS.md. The WORKSPACE.md
manifest is the only thing that knows about all three at once.

\textbf{Step 3. \texttt{cd} into each repo and run its own GSD flow.}

Inside the workspace, each repo behaves like a normal GSD project.
You change into one repo, run the normal commands, and the normal
files land in that repo's \texttt{.planning/}:

\begin{gsdbash}
cd ~/gsd-workspaces/v2-launch/api
/gsd:new-project
/gsd:discuss-phase 1
/gsd:plan-phase 1
\end{gsdbash}

\noindent
Then the same in \texttt{web-frontend} and \texttt{admin-dashboard},
each with its own plan. There is deliberately no shared planning
file across the three, because the teams work in different languages,
different release cadences, and different levels of detail. What
they share is the workspace manifest, which is updated by the
orchestrator --- never by the individual executors --- as each
repo's phase status changes.

\textbf{Step 4. Execute phases in each repo.}

Phase work happens locally in each repo. When you run
\texttt{/gsd:execute-phase} inside \texttt{api}, the executor operates
on \texttt{api}'s files and writes to \texttt{api}'s
\texttt{.planning/STATE.md}. When the phase finishes, the orchestrator
reads the final state and updates WORKSPACE.md at the workspace root
so that the cross-repo view stays accurate. If you do nothing, the
manifest will drift behind; if you run the normal GSD flow, it stays
current without any extra commands from you.

\textbf{Step 5. Strategy comparison: worktree vs. clone.}

The \texttt{--strategy} flag at workspace-creation time is worth
understanding, because the two strategies have meaningfully different
cost profiles:

\begin{center}
\small
\begin{tabularx}{\linewidth}{@{}lXXXX@{}}
\toprule
\textbf{Strategy} & \textbf{Speed} & \textbf{Disk} & \textbf{Isolation} & \textbf{Best for} \\
\midrule
worktree &
Fast (seconds) &
Minimal (shared objects) &
Shared git history &
Coordinated multi-repo planning; most launch cohorts \\
\addlinespace
clone &
Slower (full clone per repo) &
Full working copy each &
Complete &
Spikes you may throw away; same-branch parallel work \\
\bottomrule
\end{tabularx}
\end{center}

\noindent
The worktree strategy is the default because it's almost always
what you want. It shares git object storage with the parent checkout
on disk, so materializing three worktrees of large repos is
essentially free. The one constraint to be aware of is that git
refuses to have two worktrees of the same repo on the same branch
at the same time; if the workspace and your main checkout both want
to live on \texttt{main}, the workspace's worktree has to be on a
different branch.

Use the \texttt{clone} strategy when you specifically need full
isolation --- for example, when you're running a disposable
experiment and don't want git objects from the spike leaking back
into your main checkout, or when you need the workspace and main
checkout to sit on exactly the same branch at the same time.

\textbf{Step 6. List and inspect.}

When you want to see all your workspaces, or pull up the manifest
for one:

\begin{gsdbash}
/gsd:list-workspaces
cat ~/gsd-workspaces/v2-launch/WORKSPACE.md
\end{gsdbash}

\noindent
\texttt{/gsd:list-workspaces} prints a table of every workspace on
the machine, showing name, repos, strategy, and last activity.
\texttt{WORKSPACE.md} is a plain markdown file, human-readable,
that captures the orchestrator's cross-repo view --- phase status
per repo, outstanding blockers, last update timestamp. You can
skim it in any editor without running GSD.

\textbf{Step 7. Cleanup.}

When v2 finally ships, you retire the workspace:

\begin{gsdbash}
/gsd:remove-workspace v2-launch
\end{gsdbash}

\noindent
This command asks you to confirm the workspace name before
proceeding --- a small guardrail against fat-fingering the wrong
name and destroying an active workspace. Once confirmed, GSD
removes the workspace directory, tears down the worktrees (git
reports each one as pruned), and cleans up the WORKSPACE.md
manifest. The underlying repos themselves are untouched: the api,
web-frontend, and admin-dashboard checkouts elsewhere on your
disk are completely unaffected.

\begin{gsdwarning}[Workspaces refuse removal with uncommitted changes]
\texttt{/gsd:remove-workspace} will \textbf{refuse} to remove the
workspace if any constituent repo has uncommitted changes. The
protection is intentional: uncommitted work in a worktree is
easily lost if the worktree is removed before being committed
back, and ``I just meant to clean up'' is exactly how people
lose an afternoon's work. Either commit, stash, or push your
work in each repo first, then re-run the removal. If you are
absolutely sure you want to discard the uncommitted state, do
it explicitly --- stash or reset each repo by hand --- so that
the decision is deliberate.
\end{gsdwarning}

\end{gsdscenario}

\begin{gsdnote}[MCP scoping composes with workspaces]
Part II introduced project-scoped \texttt{.mcp.json} as the right
way to configure an MCP server. That pattern composes cleanly with
workspaces: a Stitch MCP server configured inside
\texttt{\~{}/gsd-workspaces/v2-launch/web-frontend/.mcp.json}
lives inside that workspace's copy of the web-frontend checkout
and does not bleed into a separate workspace's \texttt{web-frontend}
checkout, or into your main \texttt{web-frontend} checkout elsewhere
on disk. Multi-repo isolation and MCP scoping compose exactly the
way you'd expect them to.
\end{gsdnote}

\chapter{Worktree Isolation Deep Dive}

This chapter is the most technically delicate in Part III. The
details matter, and one of them in particular --- the STATE.md
ownership rule --- is the kind of thing that will silently corrupt
your planning files if you get it wrong. Read carefully.

\section{What \texttt{workflow.use\_worktrees: true} Does}

In \texttt{.planning/config.json}, a workflow setting controls how
parallel execution is sequenced:

\begin{gsdjson}
{
  "workflow": {
    "use_worktrees": true,
    "cleanup_worktrees": true
  }
}
\end{gsdjson}

\noindent
When \texttt{use\_worktrees} is \texttt{true},
\texttt{/gsd:execute-phase} creates a \emph{temporary} git worktree
for each parallel executor agent it spawns. Each agent gets its
own private filesystem view of the same git history, rooted at
a scratch directory under \texttt{.planning/worktrees/} (or wherever
the config points). The agents do their work inside those scratch
directories, and once the phase completes, the orchestrator walks
the set and aggregates their results. If \texttt{cleanup\_worktrees}
is also \texttt{true}, the scratch directories are removed after
aggregation.

Without worktrees, parallel executors would share the main checkout.
That's not a theoretical problem --- it's a concrete, reproducible
one. Two agents editing the same file race. Two agents staging
different versions of a change overwrite each other. A half-applied
refactor gets committed on top of an unrelated agent's fix and the
git history becomes unreadable. Worktrees solve this by giving each
agent a filesystem it doesn't have to share with anyone else, while
still keeping all of them rooted in the same git repository so the
orchestrator can collect their output.

\section{The STATE.md Ownership Rule}

This is the critical constraint. Get it right, and parallel execution
is safe. Get it wrong, and you reintroduce the lost-update bug that
Issue \#1571 was filed to fix.

\begin{gsdimportant}[Single-writer model for shared planning state]
The orchestrator owns STATE.md and ROADMAP.md. Parallel executor
agents must \textbf{not} write to these shared planning files
directly. Each agent writes a plan-local \textbf{SUMMARY.md} inside
its own worktree. After each wave, the orchestrator reads the
SUMMARY.md files from every worktree and aggregates them into
STATE.md and ROADMAP.md in the main checkout. This single-writer
model is the entire reason parallel GSD execution is safe. Removing
it --- even ``just this once, for a small edit'' --- reintroduces
the lost-update bugs that corrupted shared state before the fix
in Issue \#1571 landed.
\end{gsdimportant}

The rule is small enough to state in one sentence: \emph{only the
orchestrator writes to STATE.md or ROADMAP.md}. Everything else in
the design follows from that constraint.

\section{What Agents Actually Write: SUMMARY.md}

Each parallel executor agent writes exactly one file in its worktree
that is visible to the rest of the system: a plan-local SUMMARY.md.
The SUMMARY.md format is intentionally small and additive. It
captures:

\begin{itemize}
\item what the agent was asked to do (copied from its slice of the plan),
\item what it actually did (files touched, tests run, tests passing),
\item what it didn't do and why (blocking dependency, unclear requirement, environmental failure),
\item any deviations from the plan that need to be reconciled at aggregation time.
\end{itemize}

The SUMMARY.md is written in the worktree, not in the main checkout.
It lives next to the agent's work, and it describes what that agent
did within its own sandbox. The orchestrator is responsible for
reading it and deciding what to do about it. The agent never tries
to speak on behalf of the plan as a whole, and the agent never
writes into STATE.md or ROADMAP.md in the main checkout or in its
own worktree. If you find yourself writing an executor that edits
STATE.md, stop. You are reintroducing \#1571.

\section{Post-Wave Aggregation}

After all parallel executors in a wave finish, the orchestrator
runs an aggregation step. It:

\begin{enumerate}
\item Enumerates the worktrees created for this wave.
\item Reads each worktree's SUMMARY.md.
\item Reconciles the set against the master plan --- which slices
  completed, which slices deviated, which slices blocked.
\item Updates STATE.md and ROADMAP.md in the main checkout with
  the aggregated result. This is the \emph{single} write to the
  shared planning state for the entire wave.
\item If \texttt{cleanup\_worktrees} is \texttt{true}, removes the
  worktrees and prunes the git worktree registry.
\end{enumerate}

\noindent
The important property is that STATE.md goes from pre-wave to
post-wave in exactly one transition, authored by one process (the
orchestrator), reflecting the aggregate of what all the executors
did. There is never a window where two writers are touching STATE.md
at the same time, because there is only ever one writer.

\section{The \texttt{git clean} Prohibition}

\begin{gsdwarning}[\texttt{git clean} is prohibited in worktree context]
As of GSD v1.32 (Issue \#2075), the executor wrapper \textbf{blocks}
\texttt{git clean} in worktree context. Do not try to work around
this. A stray \texttt{git clean -fd} can delete the prior wave's
output before the orchestrator has had a chance to aggregate it ---
including SUMMARY.md files that were the only record of what just
happened. If you need to remove untracked files inside a worktree,
do it explicitly per-file (\texttt{rm path/to/file}), not en masse.
The prohibition is there because the failure mode is silent: you
lose work and don't notice until the aggregation step asks for a
SUMMARY.md that no longer exists on disk.
\end{gsdwarning}

The prohibition applies to any variant of the command ---
\texttt{git clean -f}, \texttt{git clean -fd}, \texttt{git clean
-fdx} --- and the wrapper will reject the call with a diagnostic
that explains why. If you have a legitimate need to remove large
quantities of untracked files from a worktree, the right move is to
finish the wave, let the orchestrator aggregate, verify that STATE.md
is up to date, and \emph{then} clean up outside of the worktree
mechanism.

\section{Toggling Worktrees Off}

Sometimes you don't want worktree-based parallel execution. Setting
\texttt{workflow.use\_worktrees} to \texttt{false} runs every
executor in the main checkout, serially. Use this when:

\begin{itemize}
\item Your environment doesn't support git worktrees cleanly. Some
  CI runners and some older Windows setups have limitations that
  make worktrees flaky. Serial execution is more boring but more
  predictable in those environments.
\item You're debugging and want a single deterministic file tree.
  Worktrees make it harder to \texttt{cd} into ``the'' checkout
  and poke around, because there are several checkouts. Turning
  them off is a reasonable diagnostic step.
\item You intentionally want sequential execution --- for example,
  on a very small machine where the memory cost of running three
  agents in parallel is worse than the wall-clock cost of running
  them in series.
\end{itemize}

The trade-off is loss of parallelism. Phases will take longer to
execute. The STATE.md ownership rule still applies in serial mode
--- the orchestrator is still the only writer --- but the risk of
races in the main checkout goes away because there is never more
than one executor running at a time. If you find yourself fighting
worktree edge cases and you don't specifically need parallelism,
the boring serial path is a legitimate choice and not a defeat.

\begin{gsdtip}[Orchestrator-only-writes is the rule]
The orchestrator-only-writes-STATE.md rule is what makes parallel
GSD execution safe. If you ever find yourself tempted to have an
executor agent write directly to STATE.md or ROADMAP.md ---
``just for this one case, to update a small field'' --- stop and
reread Issue \#1571. You are about to reintroduce the bug that
the current design exists to prevent.
\end{gsdtip}

\chapter{Team Patterns for Multi-Repo GSD}

The previous chapters were about mechanics. This chapter is about
the boring operational discipline that turns the mechanics into
something a team of humans can actually live with. None of it is
exotic. All of it is the kind of thing that quietly separates a
team with a consistent GSD workflow from a team where every repo
behaves slightly differently and no one can remember why.

\section{Standardizing Config Across Repos}

The single biggest lever for consistency is a shared
\texttt{.planning/config.json} template that every repo in your team
copies in. Without a template, every repo ends up with a slightly
different configuration --- one has worktrees enabled, another
doesn't, a third uses a different \texttt{project\_code} prefix,
a fourth has three agent skill paths the others don't know about.
The teams don't notice at first, and then six months later somebody
asks ``why does the billing repo's GSD output look different from
the accounts repo?'' and the answer is a sprawl of micro-decisions
no one wrote down.

The fix is to commit a canonical template to a shared
\texttt{gsd-templates} repo and pull it into every new project
via a setup script. The template should pin the fields that matter
for team consistency:

\begin{itemize}
\item \texttt{workflow.use\_worktrees} --- typically \texttt{true}
  for teams that execute phases in parallel.
\item \texttt{workflow.cleanup\_worktrees} --- typically \texttt{true}
  to avoid accumulating stale worktree directories.
\item \texttt{project\_code} --- the short prefix used for phase
  directories. Leave this as a template variable so each repo
  substitutes its own code at setup time.
\item Any shared agent-skill paths --- if your team has a library
  of skills living in a shared directory, the paths belong here
  so every repo picks them up.
\item Any shared verifier or reviewer configuration that your team
  applies uniformly.
\end{itemize}

What should \emph{not} go in the team template is anything specific
to a single repo: repo-specific build commands, local credentials,
one-off overrides. Those belong in the repo's own config, layered
over the template.

\section{MCP Server Scoping, Restated}

Chapter 12 touched on this in passing and it's worth restating
because it's the single most common mistake teams make with MCP in
a multi-repo context: do not try to maintain a single global
\texttt{\~{}/.claude.json} for a microservice fleet. Use
project-scoped \texttt{.mcp.json} files, one per repo, and commit
the non-sensitive parts.

The reason is that microservices are each their own context. The
billing service wants a connection to the billing database. The
auth service wants a connection to the auth database. If you wire
both of those into a single global MCP config, every Claude Code
session in every repo gets every database, and the failure modes
are confusing --- wrong DB picked up, permissions errors from a
service that should never have been able to reach that DB in the
first place, credentials leaking across contexts. Project-scoped
\texttt{.mcp.json} files keep each service's tool surface minimal
and make it obvious from the repo alone what external systems
that service talks to.

Workspaces, as we saw in Chapter 12, compose cleanly with this
pattern. Each repo inside a workspace carries its own
\texttt{.mcp.json}, and a workspace doesn't reach across its
constituent repos to share MCP servers.

\section{Workstream Naming Conventions}

Workstream names end up in \texttt{/gsd:workstreams list} output,
in \texttt{.planning/workstreams/} directory listings, in commit
messages, and occasionally in PR descriptions. A little discipline
here pays off indefinitely:

\begin{itemize}
\item \textbf{Verb-noun format.} \texttt{add-billing},
  \texttt{migrate-auth}, \texttt{refactor-orm}. Every workstream
  is \emph{doing} something, and the name should make that obvious
  at a glance.
\item \textbf{No timestamps in the name.} A workstream called
  \texttt{2026-04-backend-work} looks tidy at creation time and
  looks dreadful six months later when you have fifteen of them
  and none of the names tell you what the work was. Use creation
  timestamps via the \texttt{last touched} column in
  \texttt{/gsd:workstreams list}, not in the name.
\item \textbf{Don't use a bare ticket ID as the whole name.} A
  workstream called \texttt{JIRA-1234} is opaque. Pair the ticket
  ID with a short description: \texttt{JIRA-1234-billing-flow}.
  You get the ticket traceability and the at-a-glance readability
  at the same time, and the cost is a handful of extra characters.
\item \textbf{Short is better than long.} Aim for workstream names
  that fit in a reasonable table column. 30 characters is plenty;
  60 is painful.
\end{itemize}

None of these rules are enforced by the tooling. They are the kind
of thing you write down in your team's \texttt{gsd-templates} repo
next to the config template and agree to honor.

\section{The \texttt{project\_code} Field}

A small but underrated piece of team config is the
\texttt{project\_code} field in \texttt{.planning/config.json}.
It's a short prefix --- three or four letters --- that GSD uses
when it names phase directories and in some default commit message
templates. Setting it once:

\begin{gsdjson}
{
  "project_code": "BCS"
}
\end{gsdjson}

\noindent
means that phase directories come out looking like
\texttt{BCS-phase-1-auth-refactor} instead of just
\texttt{phase-1-auth-refactor}. When a human later skims a flat
directory listing --- or reads a commit message that references a
phase --- the prefix tells them immediately which project the
phase belongs to. In a single-repo workflow the prefix is mostly
decorative. In a multi-repo team workflow, where a reviewer might
be looking at phases from several different services in the same
week, the prefix is a small but constant cue that prevents
cross-project confusion.

Pick short, memorable codes: \texttt{BCS} for billing-core-service,
\texttt{PAY} for payments, \texttt{ADM} for admin-dashboard. Three
or four characters, upper-case, no punctuation. Record the mapping
somewhere humans can find it --- a short table in your team README
or in the \texttt{gsd-templates} repo works fine.

\section{Ticket-Based Phase Identifiers}

GSD supports tagging phases with external ticket identifiers ---
the Jira issue, the Linear ticket, the GitHub issue number, whatever
your team tracks work in. The mechanism is to record the ticket
reference as phase metadata at planning time, and GSD will carry it
through to commit message templates and PR descriptions automatically.

The details of which config key holds the ticket reference are
version-specific and worth checking against your GSD release's
documentation before you commit to a convention. What matters at
the team level is the \emph{pattern}: choose one ticket system,
pick one field name, put it in the team config template, and make
sure every planner is using it. When it's set up, you get traceable
commits without any extra effort from the people doing the work ---
the ticket ID threads through from plan to commit to PR automatically.

The payoff shows up when you're looking at a six-month-old
commit and wondering why it exists. If the commit says ``implement
session token rotation'' and the message also carries
\texttt{PAY-1734}, you click through to the ticket and get the full
context. Without that thread, you're archaeologizing.

\begin{gsdsuccess}[Standardize the boring parts]
Standardize the boring parts --- config template, naming
conventions, MCP scoping, ticket integration --- and your team's
GSD output will look consistent across repos even when different
humans are driving different services. The interesting parts ---
what to plan, how to scope phases, when to escalate a workstream
to a workspace --- stay where they belong: in the hands of the
people doing the work. Operational consistency and creative
autonomy are not in tension. They live at different layers of the
stack, and a well-set-up team separates them cleanly.
\end{gsdsuccess}
% ============================================================================
%  appendices.tex
%  Three quick-reference appendices for "GSD: Going Further".
% ============================================================================

\appendix

\chapter{Runtime Compatibility Matrix}
\label{app:runtime-matrix}

A consolidated lookup of every coding-agent runtime that GSD's installer
targets. Use this when you need to verify the install path, the verify
command, or the workspace/MCP support level for a runtime without paging
back through Part~I.

\begin{small}
\begin{longtable}{@{}p{2.4cm}p{2.4cm}p{3.6cm}p{1.8cm}p{1.4cm}p{1.4cm}@{}}
\toprule
\textbf{Runtime} & \textbf{Skills Format} & \textbf{Install Dir (global)} &
\textbf{Verify} & \textbf{MCP} & \textbf{Workspace} \\
\midrule
\endfirsthead
\toprule
\textbf{Runtime} & \textbf{Skills Format} & \textbf{Install Dir (global)} &
\textbf{Verify} & \textbf{MCP} & \textbf{Workspace} \\
\midrule
\endhead
Claude Code 2.1.88+ & \texttt{skills/gsd-*/SKILL.md} &
\texttt{\textasciitilde/.claude/skills/} & \texttt{/gsd:help} & yes & yes \\
Claude Code (legacy) & \texttt{commands/gsd/*.md} &
\texttt{\textasciitilde/.claude/commands/} & \texttt{/gsd:help} & yes & yes \\
Codex & \texttt{skills/gsd-*/SKILL.md} &
\texttt{\textasciitilde/.codex/skills/} & \texttt{\$gsd-help} & partial & yes \\
Copilot & prompts + agents &
\texttt{\textasciitilde/.github/} & \texttt{/gsd:help} & partial & no \\
Cursor & \texttt{skills/gsd-*/SKILL.md} &
\texttt{\textasciitilde/.cursor/} & \texttt{/gsd:help} & yes & yes \\
Windsurf & markdown transform &
\texttt{\textasciitilde/.windsurf/} & \texttt{/gsd:help} & partial & yes \\
OpenCode & config file &
\texttt{\textasciitilde/.config/opencode/} & \texttt{/gsd-help} & yes & yes \\
Gemini CLI & skills &
\texttt{\textasciitilde/.gemini/} & \texttt{/gsd:help} & yes & yes \\
Antigravity & skills &
\texttt{\textasciitilde/.gemini/antigravity/} & \texttt{/gsd:help} & yes & yes \\
Cline & \texttt{.clinerules} &
\texttt{\textasciitilde/.cline/} & auto-loaded & no & partial \\
Augment & skills transform &
\texttt{\textasciitilde/.augment/} & \texttt{/gsd:help} & partial & yes \\
Qwen Code & skills (open standard) &
\texttt{\textasciitilde/.qwen/skills/} & \texttt{/gsd:help} & yes & yes \\
\bottomrule
\end{longtable}
\end{small}

\textbf{Notes.} The \texttt{MCP} column reflects whether the host runtime
exposes Model Context Protocol tool surfaces to its agents. \texttt{partial}
means MCP is supported via configuration but the host's discovery story is
limited (e.g., Cline reads its rules file but does not expose MCP tools the
same way Claude Code does). The \texttt{Workspace} column tracks GSD
multi-repo workspace support; \texttt{partial} means the runtime accepts
GSD's installed skill files but does not run worktree-based phases.

\begin{gsdtip}
Antigravity has both a global install directory and a per-project
\texttt{./.agent/} directory. GSD's installer auto-detects which to use
based on whether you run the install script with \texttt{--global} or from
inside a project directory.
\end{gsdtip}

\chapter{MCP Server Quick-Reference}
\label{app:mcp-reference}

A scan-friendly table of MCP servers used in this guide, plus a couple of
honourable mentions you'll reach for early. Each row lists the package, the
transport, the auth model, and the minimal config snippet you'd drop into
\texttt{.mcp.json}.

\begin{small}
\begin{tabularx}{\textwidth}{@{}llllX@{}}
\toprule
\textbf{Server} & \textbf{Package} & \textbf{Transport} & \textbf{Auth} &
\textbf{Minimal Config Snippet} \\
\midrule
Google Stitch &
\texttt{@\_davideast/stitch-mcp} &
stdio (proxy) &
OAuth / API key &
\texttt{\{"command":"npx","args":["@\_davideast/stitch-mcp","proxy"]\}} \\
\addlinespace
Postgres &
\texttt{@modelcontextprotocol/} \texttt{server-postgres} &
stdio &
connection string &
\texttt{\{"command":"npx","args":["-y","@modelcontextprotocol/server-postgres"],"env":\{"DATABASE\_URL":"\$\{DATABASE\_URL\}"\}\}} \\
\addlinespace
Filesystem &
\texttt{@modelcontextprotocol/} \texttt{server-filesystem} &
stdio &
none (path scoped) &
\texttt{\{"command":"npx","args":["-y","@modelcontextprotocol/server-filesystem","/abs/path"]\}} \\
\addlinespace
GitHub &
\texttt{@modelcontextprotocol/} \texttt{server-github} &
stdio &
PAT &
\texttt{\{"command":"npx","args":["-y","@modelcontextprotocol/server-github"],"env":\{"GITHUB\_PERSONAL\_ACCESS\_TOKEN":"\$\{GITHUB\_TOKEN\}"\}\}} \\
\addlinespace
Brave Search &
\texttt{@modelcontextprotocol/} \texttt{server-brave-search} &
stdio &
API key &
\texttt{\{"command":"npx","args":["-y","@modelcontextprotocol/server-brave-search"],"env":\{"BRAVE\_API\_KEY":"\$\{BRAVE\_API\_KEY\}"\}\}} \\
\bottomrule
\end{tabularx}
\end{small}

\begin{gsdimportant}
Every snippet that contains a credential reference here uses an environment
variable placeholder. Drop the snippet into a real \texttt{.mcp.json} only
after exporting the corresponding variable in your shell, and add
\texttt{.mcp.json} to \texttt{.gitignore} if it carries any credentials.
\end{gsdimportant}

\textbf{Choosing a transport.} stdio is the right default for tools that
live next to your code. Use HTTP transport (\texttt{--transport http}) for
cloud services and shared team servers; the GSD User Guide and the project's
\texttt{docs/USER-GUIDE.md} cover the HTTP setup in more depth.

\chapter{Workspace Command Cheat Sheet}
\label{app:cheatsheet}

The dense version. Print this page and tape it next to your monitor.

\section*{Workstreams (single-repo, isolated planning state)}

\begin{gsdplain}
/gsd:workstreams create <name>      Create a new workstream
/gsd:workstreams switch <name>      Switch the active workstream
/gsd:workstreams list               Show all workstreams (active + archived)
/gsd:workstreams complete <name>    Archive a workstream when done
\end{gsdplain}

State lives under \texttt{.planning/workstreams/<name>/}. Switching is fast --
GSD just swaps which subtree it reads.

\section*{Workspaces (multi-repo or full-isolation)}

\begin{gsdplain}
/gsd:new-workspace --name <name> --repos <repo1,repo2,...>
                       [--strategy worktree|clone]
/gsd:list-workspaces
/gsd:remove-workspace <name>
\end{gsdplain}

Strategies:
\begin{itemize}[leftmargin=2em,itemsep=0.2em]
  \item \texttt{worktree} (default) -- git worktrees, lightweight, shares git
        history. Cannot have two worktrees on the same branch.
  \item \texttt{clone} -- full \texttt{git clone}, complete isolation, higher
        disk cost. Use for spikes and throw-aways.
\end{itemize}

Workspaces live under \texttt{\textasciitilde/gsd-workspaces/<name>/} with a
\texttt{WORKSPACE.md} manifest at the root.

\begin{gsdwarning}
\texttt{/gsd:remove-workspace} refuses removal if any constituent repo has
uncommitted changes. Commit, stash, or push first.
\end{gsdwarning}

\section*{Relevant config keys}

\begin{gsdplain}
.planning/config.json:

  workflow.use_worktrees: true | false
      Per-phase parallel executor isolation via git worktrees.

  workflow.cleanup_worktrees: true | false
      Whether the orchestrator removes worktrees after phase completion.

  project_code: "<prefix>"
      Short prefix used in phase directory names and commit messages.

  agent_skills:
    executor:   ["./skills/<dir>/", ...]
    planner:    ["./skills/<dir>/", ...]
    researcher: ["./skills/<dir>/", ...]
      Inject custom instruction directories into specific subagent types.
\end{gsdplain}

\section*{The single rule that makes parallel execution safe}

\begin{gsdimportant}
The orchestrator owns STATE.md and ROADMAP.md. Parallel executor agents
write only their plan-local SUMMARY.md. The orchestrator aggregates after
each wave. Issue \#1571 documents what happens when this rule is broken --
don't be the one who reintroduces it.
\end{gsdimportant}

\end{document}
